\documentclass[conference]{IEEEtran}
\IEEEoverridecommandlockouts
% The preceding line is only needed to identify funding in the first footnote. If that is unneeded, please comment it out.
%Template version as of 6/27/2024



\def\BibTeX{{\rm B\kern-.05em{\sc i\kern-.025em b}\kern-.08em
    T\kern-.1667em\lower.7ex\hbox{E}\kern-.125emX}}

\usepackage{amsmath,amssymb,amsfonts}
\usepackage{algorithmic}
\usepackage{mathtools}
\usepackage{algorithm}
\usepackage{array}
\usepackage[font=normalsize]{subfig}
\usepackage{textcomp}
\usepackage{stfloats}
\usepackage{url}
\usepackage{verbatim}
\usepackage{siunitx}
\usepackage{graphicx}
\usepackage{cite}
\usepackage{makecell}
\usepackage[switch]{lineno} % easier for the review line numbers
\usepackage{booktabs}
\newcommand{\ra}[1]{\renewcommand{\arraystretch}{#1}}
\hyphenation{op-tical net-works semi-conduc-tor IEEE-Xplore}
\usepackage[dvipsnames]{xcolor}
\usepackage{placeins}
\usepackage{xifthen}
\usepackage{colortbl}
\usepackage{bbm}
\usepackage{array}
\newcolumntype{H}{>{\setbox0=\hbox\bgroup}c<{\egroup}@{}}
\newcommand{\prob}[1][]{%requires xifthen package%
\ifthenelse{\isempty{#1}}%
      {\ensuremath{P}}%
    {\ensuremath{P\left\(#1\right\)}}%
}
\newcommand{\vect}[1]{\boldsymbol{\mathrm{#1}}}
\newcommand{\mat}[1]{\boldsymbol{\mathrm{#1}}}
\newcommand{\tens}[1]{\boldsymbol{\mathsf{{#1}}}}\newcommand{\MSE}{\mathrm{MSE}}
\newcommand{\moddef}{\mathrm{mod}}
\newcommand{\vecop}{\text{vec}}
\newcommand{\CP}{L}
\newcommand{\hddots}{\hdots}
\newcommand{\tr}[1]{\mathrm{tr}\left(#1\right)}
\newcommand{\diag}[1]{\mathrm{diag}\left(#1\right)}
\newcommand*{\inC}[1]{\in\mathbb{C}^{#1}}
\newcommand{\norm}[1]{\left\lVert#1\right\rVert}
\newcommand{\abs}[1]{\left\lvert#1\right\rvert}
\newcommand{\expt}[1]{\mathbb{E}\left(#1\right)}
\newcommand{\cn}[2]{\ensuremath{\mathcal{C}\mathcal{N}\left(#1,#2\right)}}
\newcommand{\cw}[3]{\ensuremath{\mathcal{C}\mathcal{W}\left(#1,#2,#3\right)}}
\newcommand{\icw}[3]{\ensuremath{{\mathcal{C}\mathcal{W}}^{-1}\left(#1,#2,#3\right)}}
\DeclareMathOperator*{\maximize}{maximize}
\DeclareMathOperator*{\minimize}{minimize}
\DeclareMathOperator*{\argmax}{arg\,max}
\newcommand{\todo}[1]{{\color{red} TODO {#1}}}
\newcommand{\thomas}[1]{{\color{blue}[Thomas: #1]}} 
\newcommand{\review}[1]{{{#1}}}
\newcommand{\gilles}[1]{{\color{BurntOrange}[Gilles: #1]}} 

\usepackage{pgfplots}
\usepackage{nicefrac}
\usepgfplotslibrary{groupplots,dateplot}
%\usetikzlibrary{snakes}

\pgfplotsset{compat=1.18}
\usepgfplotslibrary{fillbetween}
\usepackage{tikz}
\usetikzlibrary{positioning}
\usepackage{hyperref}
\usepackage{acronym}
\usepackage[capitalize]{cleveref}
% \usepackage{subfig}
% \captionsetup[subfloat]{labelfont=small,labelfont=rm}
\usepackage{subcaption} % put this in your preamble

\usetikzlibrary{patterns}
\usepgfplotslibrary{polar}
\DeclareMathOperator*{\argmin}{arg\,min}
\newcommand{\vertfig}[2][]{%
  \begin{minipage}{5in}\subfloat[#1]{#2}\end{minipage}}
\newcommand{\horizfig}[2][]{%
  \begin{minipage}{3in}\subfloat[#1]{#2}\end{minipage}}
\usetikzlibrary{shapes}
\usepackage{pgfplotstable}
\usepackage{filecontents}
\usepackage[acronym]{glossaries}
\newacronym[plural=BSs,firstplural=base stations (BSs)]{bs}{BS}{base station}
\newacronym[plural=PAs,firstplural=power amplifiers (PAs)]{pa}{PA}{power amplifier}
\newacronym[plural=APs,firstplural=access points (APs)]{ap}{AP}{access point}
\newacronym{sinr}{SINR}{signal-to-interference-plus-noise ratio}
\newacronym{csi}{CSI}{channel state information}
\newacronym{sdr}{SDR}{signal-to-distortion ratio}
\newacronym{snr}{SNR}{signal-to-noise ratio}
\newacronym{sndr}{SNDR}{signal-to-noise-and-distortion ratio}
\newacronym{snidr}{SNIDR}{signal-to-noise-and-interference-and-distortion ratio}
\newacronym{los}{LOS}{line-of-sight}
\newacronym{z3ro}{Z3RO}{zero third-order distortion}
\newacronym{mimo}{MIMO}{multiple input multiple output}
\newacronym{mrt}{MRT}{maximum ratio transmission}
\newacronym{dpd}{DPD}{digital pre-distortion}

\newacronym{ula}{ULA}{uniform linear array}

\newacronym{psd}{PSD}{power spectral density}
\newacronym{aclr}{ACLR}{adjacent channel leakage ratio}
\newacronym{cpu}{CPU}{central processing unit}
\newacronym{iid}{i.i.d.}{independently and identically distributed}
\newacronym{ue}{UE}{user equipment}
\newacronym{nlos}{NLoS}{non-line-of-sight}
\newacronym{hpbm}{HPBM}{half power beam width}
\newacronym{rts}{RTS}{ray tracing simulator}
\newacronym{ag}{AG}{array gain}
\newacronym{arp}{ARP}{Antenna Reference Point}
\newacronym{ecdf}{eCDF}{empirical cumulative distribution function}
\newacronym{evm}{EVM}{Error Vector Magnitude}
\newacronym{fft}{FFT}{fast Fourier transform}
\newacronym{ib}{IB}{in-band}
\newacronym{im}{IM}{intermodulation}
\newacronym{mf}{MF}{matched filterifng}
\newacronym{mr}{MR}{maximum ratio}
\newacronym{ofdm}{OFDM}{orthogonal frequency division multiplexing}
\newacronym{oob}{OOB}{out-of-band}
\newacronym{papr}{PAPR}{peak-to-average power ratio}
\newacronym{rx}{RX}{Receiver}
\newacronym{trp}{TRP}{Transmission Reception Point}
\newacronym{tx}{TX}{transmitter}
\newacronym{zf}{ZF}{zero forcing}
\newacronym{dl}{DL}{downlink}
\newacronym{rzf}{RZF}{regularized zero forcing}
\newacronym{slc}{SLC}{spatial leakage suppression}
\newacronym{kpi}{KPI}{key performance indicator}
\newacronym{awgn}{AWGN}{additive white Gaussian noise}
\newacronym{qam}{QAM}{quadrature amplitude modulation}
\newacronym{amam}{AM/AM}{amplitude modulation to amplitude modulation}
\newacronym{ampm}{AM/PM}{amplitude modulation to phase modulation}
\newacronym{zmcscg}{ZMCSCG}{zero mean circularly symmetric complex Gaussian}

\newacronym{gnn}{GNN}{graph neural network}
\newacronym{ccnn}{CCNN}{circular convolutional neural network}
\newacronym{lrelu}{LReLU}{leaky rectified linear unit }
\newacronym{sdg}{SDG}{Sustainable Development Goal}
\newacronym{ibo}{IBO}{input back-off}
\newacronym{cdf}{CDF}{cumulative distribution function}
\newacronym{nn}{NN}{neural network}
\newacronym{mlp}{MLP}{multilayer perceptron}
\newacronym{mmimo}{mMIMO}{massive MIMO}
\newacronym{dab}{DAB}{distortion-aware beamforming}
\newacronym{flops}{FLOPs}{floating point operations}
\newacronym{mmwave}{mmWave}{millimeter-wave}
\newacronym{gops}{GFLOPs/s}{giga floating point operations per second}
\newacronym{sp}{SP}{single precision}
\newacronym{dsp}{DSP}{digital signal processing}
\newacronym{dac}{DAC}{digital-to-analog converter}
\newacronym{aqnm}{AQNM}{additive quantization noise model}
\newacronym{sinqdr}{SINQDR}{signal-to-interference-noise-and-quantization-distortion ratio}
\newacronym{nmsqe}{NMSQE}{normalized mean squared quantization error}
\newacronym{mmsqe}{MMSQE}{minimum mean squared quantization error}
\newacronym{agc}{AGC}{automatic gain control}
\newacronym{rf}{RF}{radio frequency}
\newacronym{sfdr}{SFDR}{spurious free dynamic range}
\newacronym{adc}{ADC}{analog-to-digital converter}
\newacronym{mmse}{MMSE}{minimum mean squared error}
\newacronym{cnn}{CNN}{convolutional neural network}
\newacronym{mse}{MSE}{mean squared error}
\newacronym{nmse}{NMSE}{normalized mean squared error}
\newacronym{nco}{NCO}{neural combinatorial optimization}
\newacronym{rfdac}{RF-DAC}{radio frequency digital-to-analog converter}
\newacronym{idft}{IDFT}{inverse discrete Fourier transform}
\newacronym{dft}{DFT}{discrete Fourier transform}
\newacronym{pdp}{PDP}{power delay profile}
\newacronym{tdl}{TDL}{tapped delay line}
\newacronym{tdd}{TDD}{time division duplexing}
\glsdisablehyper
\usepackage{amsthm}
\newtheoremstyle{mystyle}%                % Name
  {}%                                     % Space above
  {}%                                     % Space below
  {\itshape}%                             % Body font
  {}%                                     % Indent amount
  {\bfseries}%                            % Theorem head font
  {.}%                                    % Punctuation after theorem head
  { }%                                    % Space after theorem head, ' ', or \newline
  {}%                                     % Theorem head spec (can be left empty, meaning `normal')
\theoremstyle{mystyle}
\newtheorem{theorem}{Theorem}
\newtheorem{lemma}{Lemma}
\newtheorem{corollary}{Corollary}
\newtheorem{proposition}{Proposition}
\newtheorem{remark}{Remark}
\newtheorem{assumption}{Assumption}
\usepackage[font=footnotesize]{caption}

% fix for edas error
%\setlength{\columnsep}{0.12in}  % 0.12 inches

    
\begin{document}

\title{Notes Power Model D-MIMO\\
%Distortion-Aware GNN Precoding for Multi-Carrier Large Antenna Systems with Adaptive Precoder Selection
%Extending Distortion-Aware Graph Neural Network Precoding to Multi-Carrier Systems with Channel-Adaptive Precoder Selection
%Learning When to Learn: GNN Distortion-Aware Precoding and Adaptive Switching for Nonlinear Multi-Carrier MIMO
%Learning When to Learn: Distortion-Aware GNN Precoding for Multi-Carrier Large Antenna Systems with Adaptive Precoder Selection

\thanks{thanks
}
}

% }
\author{\IEEEauthorblockN{Thomas Feys
	}
	\IEEEauthorblockA{KU Leuven,  Dept. ESAT-WaveCore, Dramco, B-9000 Ghent, Belgium
	}
    %\IEEEauthorblockA{E-mail: thomas.feys@kuleuven.be}
}

% \author{\IEEEauthorblockN{1\textsuperscript{st} Given Name Surname}
% \IEEEauthorblockA{\textit{dept. name of organization (of Aff.)} \\
% \textit{name of organization (of Aff.)}\\
% City, Country \\
% email address or ORCID}
% \and
% \IEEEauthorblockN{2\textsuperscript{nd} Given Name Surname}
% \IEEEauthorblockA{\textit{dept. name of organization (of Aff.)} \\
% \textit{name of organization (of Aff.)}\\
% City, Country \\
% email address or ORCID}
% \and
% \IEEEauthorblockN{3\textsuperscript{rd} Given Name Surname}
% \IEEEauthorblockA{\textit{dept. name of organization (of Aff.)} \\
% \textit{name of organization (of Aff.)}\\
% City, Country \\
% email address or ORCID}
% \and
% \IEEEauthorblockN{4\textsuperscript{th} Given Name Surname}
% \IEEEauthorblockA{\textit{dept. name of organization (of Aff.)} \\
% \textit{name of organization (of Aff.)}\\
% City, Country \\
% email address or ORCID}
% \and
% \IEEEauthorblockN{5\textsuperscript{th} Given Name Surname}
% \IEEEauthorblockA{\textit{dept. name of organization (of Aff.)} \\
% \textit{name of organization (of Aff.)}\\
% City, Country \\
% email address or ORCID}
% \and
% \IEEEauthorblockN{6\textsuperscript{th} Given Name Surname}
% \IEEEauthorblockA{\textit{dept. name of organization (of Aff.)} \\
% \textit{name of organization (of Aff.)}\\
% City, Country \\
% email address or ORCID}
% }

\maketitle

% \begin{abstract}

% \end{abstract}

% \begin{IEEEkeywords}
% \end{IEEEkeywords}

\glsresetall
\section{Central MIMO Power Model}
\label{sec:pwr_model}

We adopt the parametric power model for upper mid-band (FR3) base stations of
\cite{peschiera2026fr3}. The base station (BS) is equipped with
$M_{\mathrm{ant}}$ antennas, an equal number of power amplifiers (PAs, used in
the downlink) and low-noise amplifiers (LNAs, used in the uplink), and
$M_{\mathrm{RF}}$ radio-frequency (RF) chains. It serves $K$ single-antenna
users over a bandwidth $B$ at carrier frequency $f_c$ using OFDM and
space-division multiple access. Setting $M_{\mathrm{RF}}=M_{\mathrm{ant}}$
selects fully-digital beamforming; $M_{\mathrm{RF}}<M_{\mathrm{ant}}$ selects
hybrid partially-connected beamforming with $M_{\mathrm{ant}}/M_{\mathrm{RF}}$
phase shifters per RF chain.

\subsection{Model structure}
The total consumed power is split into three hardware blocks; digital signal
processing, analog signal processing, and the PAs, each divided by its
supply-and-cooling efficiency $\eta_{c,\mathrm{sc}}\in(0,1]$:
\begin{equation}
\label{eq:pcons}
P_{\mathrm{cons}} =
\frac{\overline{P_{\mathrm{dig}}}}{\eta_{\mathrm{dig,sc}}}
+ \frac{\overline{P_{\mathrm{ana}}}}{\eta_{\mathrm{ana,sc}}}
+ \frac{\overline{P_{\mathrm{PA}}}}{\eta_{\mathrm{PA,sc}}} .
\end{equation}
Each block is built from the working-mode (active) consumption of its
subcomponents, averaged over the operating modes within a frame. We first
define this frame averaging (\Cref{sec:frame}) and then reuse it as a
function for every component (\Cref{sec:pa,sec:ana,sec:dig}).

\subsection{Frame structure and operating-mode averaging}
\label{sec:frame}
A frame of duration $T_f$ contains a downlink (DL) and an uplink (UL) phase
under \gls{tdd}, with duration ratios $\tau_{\mathrm{DL}},\tau_{\mathrm{UL}}\in[0,1]$
and $\tau_{\mathrm{UL}}=1-\tau_{\mathrm{DL}}$. Within each direction
$i\in\{\mathrm{DL},\mathrm{UL}\}$ the time is further subdivided into three
operating modes, plus an idle mode when the direction is inactive:
\begin{itemize}
  \item \emph{reference signalling}, for a time fraction
        $\tau_i\,\tau_{i,\mathrm{sig}}$;
  \item \emph{data} transmission/reception, for a fraction
        $\tau_i\,(1-\tau_{i,\mathrm{sig}})\,\dfrac{N_{a,i}}{N_i}$;
  \item \emph{micro-sleep} (DTX), for a fraction
        $\tau_i\,(1-\tau_{i,\mathrm{sig}})\Big(1-\dfrac{N_{a,i}}{N_i}\Big)$;
  \item \emph{idle}, for a fraction $(1-\tau_i)$,
\end{itemize}
where $N_i$ is the number of OFDM symbols (time slots) of direction $i$ and
$N_{a,i}\le N_i$ the subset that are \emph{active} (carrying data). The
\emph{active-resource time ratio} $N_{a,i}/N_i$ is what sets the
data-versus-micro-sleep split: a slot is either active or not, regardless of how
many users share it.
with $\tau_{i,\mathrm{sig}}\in[0,1]$ the signalling fraction of direction $i$.

\paragraph{Average physical resource load}
The data-mode weight is more conveniently expressed through the
\emph{average physical resource load} $\bar x_i$, which we now relate to the
time ratio $N_{a,i}/N_i$. Let $K_{\max}$ be the maximum number of users that can
be served (the number of spatial streams) and $Q_{\max}$ the number of available
data subcarriers per active OFDM symbol. The instantaneous load of user $k$ in
time slot $n$ is the ratio of utilised to available data subcarriers,
\begin{equation}
\label{eq:load_inst}
\ell_i[k,n] = \frac{Q_i[k,n]}{Q_{\max}},\qquad i\in\{\mathrm{DL},\mathrm{UL}\},
\end{equation}
for $k=1,\dots,K_{\max}$ and $n=1,\dots,N_i$. Averaging over all $K_{\max}$
potential users and $N_i$ slots gives the average physical resource load
\begin{equation}
\label{eq:load_avg}
\bar x_i \triangleq \bar\ell_i
= \frac{1}{N_i}\sum_{n=1}^{N_i}\left(
\frac{1}{K_{\max}}\sum_{k=1}^{K_{\max}}\frac{Q_i[k,n]}{Q_{\max}}\right).
\end{equation}
Under space-division multiple access (no OFDMA), an active user in an active
slot is allocated \emph{all} data subcarriers, $Q_i[k,n]=Q_{\max}$. Indexing the
$K_i$ active users as $\{1,\dots,K_i\}$ and the $N_{a,i}$ active slots as
$\{1,\dots,N_{a,i}\}$, each active-user/active-slot pair contributes $1$ to
\eqref{eq:load_avg} and every other pair $0$; the double sum equals $K_i N_{a,i}$
and
\begin{equation}
\label{eq:load_simplified}
\bar x_i = \frac{N_{a,i}}{N_i}\,\frac{K_i}{K_{\max}}.
\end{equation}

\paragraph{Equivalence and modelling assumption}
The duty-cycle split in the frame structure is governed by the active-resource
time ratio $N_{a,i}/N_i$, whereas the power model uses the average load $\bar
x_i$ as the data-mode weight. Comparing \eqref{eq:load_simplified} with the data
fraction $\tau_i(1-\tau_{i,\mathrm{sig}})\,N_{a,i}/N_i$, the two agree,
\begin{equation}
\label{eq:equiv}
\bar x_i = \frac{N_{a,i}}{N_i}
\quad\Longleftrightarrow\quad
\frac{K_i}{K_{\max}} = 1
\quad\Longleftrightarrow\quad
K_i = K_{\max},
\end{equation}
i.e.\ only when every spatial stream is scheduled. When the maximum number of
users is active the average load is simply the ratio of active to total slots;
for $K_i<K_{\max}$ the average load \eqref{eq:load_simplified} is the true active
fraction scaled by $K_i/K_{\max}<1$. The number of active users $K_i$ separately
sets the transmit power and the achievable rate, but not whether a slot is
active. We use $\bar x_i$ as the data-mode weight throughout (and $1-\bar x_i$ as
the micro-sleep/DTX fraction); this is therefore \emph{exact} in the
full-spatial-load regime $K_i=K_{\max}$ evaluated here and an approximation
otherwise.

In micro-sleep and idle a component drops to a fraction of its working
consumption, set by the reduction factors
$\delta^{\mu}_c,\delta^{0}_c\in[0,1]$. For a component $c$ with working power
$P^{\mathrm a}$ during data, $P^{\mathrm s}$ during signalling, and base power
$P^{0}$ scaled in the reduced modes, the frame-averaged consumption of
direction $i$ is the time-weighted sum
\begin{align}
\label{eq:frameavg}
\mathcal{A}_i^{c}\!\left(P^{\mathrm a},P^{\mathrm s},P^{0}\right) \triangleq{}&
\bar x_i\,\tau_i\,(1-\tau_{i,\mathrm{sig}})\,P^{\mathrm a}
+ \tau_i\,\tau_{i,\mathrm{sig}}\,P^{\mathrm s} \nonumber\\
&+ (1-\bar x_i)\,\tau_i\,(1-\tau_{i,\mathrm{sig}})\,\delta^{\mu}_c\,P^{0}
\nonumber\\
&+ (1-\tau_i)\,\delta^{0}_c\,P^{0} .
\end{align}
For the digital and analog blocks the consumption is identical in all working
modes, $P^{\mathrm a}=P^{\mathrm s}=P^{0}=P$, and we use the shorthand
$\mathcal{A}_i^{c}(P)\triangleq\mathcal{A}_i^{c}(P,P,P)$. Only the PA
distinguishes the three levels (\Cref{sec:pa}).

Each direction's hardware is averaged over the \emph{whole} frame, and the two
directions are accounted for separately, e.g.\ for the digital block
$\overline{P_{\mathrm{dig}}}=\mathcal{A}_{\mathrm{DL}}^{\mathrm{dig}}(P_{\mathrm{dig,DL}})
+\mathcal{A}_{\mathrm{UL}}^{\mathrm{dig}}(P_{\mathrm{dig,UL}})$. The idle term
$(1-\tau_i)\,\delta^{0}_c\,P^{0}$ is therefore the reduced consumption of
direction $i$'s chain while that direction is \emph{not} active, i.e.\ while the
opposite direction occupies the frame. For $i=\mathrm{DL}$ this fraction equals
$\tau_{\mathrm{UL}}$ only because the gap-free TDD frame satisfies
$\tau_{\mathrm{UL}}=1-\tau_{\mathrm{DL}}$; it represents the idle DL chain, not
the active UL processing, which is captured by the separate
$\mathcal{A}_{\mathrm{UL}}$ term.

\begin{remark}[Load-dependent split]
Only the data term in \eqref{eq:frameavg} scales with the load $\bar x_i$,
while the micro-sleep term scales with $(1-\bar x_i)$. Collecting the terms
proportional to $\bar x_i$ separates each averaged consumption into a
\emph{load-dependent} part and a \emph{load-independent} part, which is used in
the analysis to attribute consumption to the delivered traffic.
\end{remark}

\subsection{Power amplifiers}
\label{sec:pa}
The PAs transmit only in the DL. With a per-PA output power
$P_a=P_{T,\mathrm{DL}}/M_{\mathrm{ant}}$, where $P_{T,\mathrm{DL}}$ is the total
DL transmit power, the instantaneous consumption of one PA is modelled as
\begin{equation}
\label{eq:pa}
P_{\mathrm{PA}}(p) = \frac{1}{\eta_{\mathrm{PA,max}}}
\left[\xi\,P_{\max} + (1-\xi)\,P_{\max}^{\,1-\alpha}\,p^{\alpha}\right],
\end{equation}
where $p$ is the output power, $P_{\max}$ the maximum output power (at a fixed
back-off from saturation), $\eta_{\mathrm{PA,max}}$ the maximum PA efficiency,
$\xi\in[0,1]$ the static-to-dynamic consumption ratio, and $\alpha\in[0,1]$
the efficiency exponent. Averaging over the DL frame and summing over the
$M_{\mathrm{ant}}$ PAs gives
\begin{equation}
\label{eq:pa_avg}
\overline{P_{\mathrm{PA}}} = M_{\mathrm{ant}}\,
\mathcal{A}_{\mathrm{DL}}^{\mathrm{PA}}\!\left(
P_{\mathrm{PA}}(P_a),\,
P_{\mathrm{PA}}(\zeta_{\mathrm{sig}}P_{\max}),\,
P_{\mathrm{PA}}(0)\right),
\end{equation}
where $\zeta_{\mathrm{sig}}\in(0,1]$ is the fraction of $P_{\max}$ used during
signalling, and the zero-output consumption $P_{\mathrm{PA}}(0)$ is the base
level scaled in micro-sleep and idle.

\subsection{Analog signal processing}
\label{sec:ana}
The analog block comprises, per RF chain, the DACs (DL) or ADCs (UL), RF
filters, mixers, $M_{\mathrm{PS}}$ phase shifters, and the LNAs (UL); the local
oscillator (LO) is shared across chains. With
$M_{\mathrm{PS}}=M_{\mathrm{ant}}/M_{\mathrm{RF}}$ for hybrid beamforming
($M_{\mathrm{PS}}=0$ for fully-digital), the working-mode consumption per
direction is
\begin{align}
\label{eq:ana_dl}
P_{\mathrm{ana,DL}} &= P_{\mathrm{LO}} + M_{\mathrm{RF}}\big(
2P_{\mathrm{DAC}} + 2P_{\mathrm{filt,RF}} + 2P_{\mathrm{mix}} \nonumber\\
&\qquad + M_{\mathrm{PS}}P_{\mathrm{PS}}\big),\\
\label{eq:ana_ul}
P_{\mathrm{ana,UL}} &= P_{\mathrm{LO}} + M_{\mathrm{RF}}\Big(
2P_{\mathrm{ADC}} + 2P_{\mathrm{filt,RF}} + 2P_{\mathrm{mix}} \nonumber\\
&\qquad + M_{\mathrm{PS}}P_{\mathrm{PS}}
+ \tfrac{M_{\mathrm{ant}}}{M_{\mathrm{RF}}}P_{\mathrm{LNA}}\Big),
\end{align}
with the factor $2$ accounting for the in-phase and quadrature branches. The
subcomponents are modelled as
\begin{align}
P_{\mathrm{DAC}} &= \Xi_{\mathrm{DAC},1}2^{b_{\mathrm{DAC}}}
+ \Xi_{\mathrm{DAC},2}\,b_{\mathrm{DAC}}f_{\mathrm{DAC}},\\
P_{\mathrm{ADC}} &= W_{\mathrm{ADC}}\,f_{\mathrm{ADC}}\,2^{b_{\mathrm{ADC}}},\\
P_{\mathrm{mix}} &= \Xi_{\mathrm{mix}}f_c,\\
P_{\mathrm{PS}} &= \Xi_{\mathrm{PS}}B,\\
P_{\mathrm{LNA}} &= \Xi_{\mathrm{LNA}}B,
\end{align}
where $b_{\mathrm{DAC}},b_{\mathrm{ADC}}$ are the effective resolutions,
$f_{\mathrm{DAC}},f_{\mathrm{ADC}}$ the converter sampling rates,
$W_{\mathrm{ADC}}$ the ADC Walden figure-of-merit, and $\Xi_\bullet$ the
respective figures-of-merit. The frame-averaged analog consumption is
\begin{equation}
\label{eq:ana_avg}
\overline{P_{\mathrm{ana}}} =
\mathcal{A}_{\mathrm{DL}}^{\mathrm{ana}}\!\left(P_{\mathrm{ana,DL}}\right)
+ \mathcal{A}_{\mathrm{UL}}^{\mathrm{ana}}\!\left(P_{\mathrm{ana,UL}}\right)
+ P_{\mathrm{ana,misc}} .
\end{equation}

\subsection{Digital signal processing}
\label{sec:dig}
In the DL the digital block runs the channel encoder, MIMO precoder, IFFT, DPD,
and baseband (BB) filter; in the UL the channel decoder, MIMO combiner, FFT,
and BB filter. Each subcomponent consumption is a static (signal-independent)
term plus a dynamic term equal to its complex giga-operations-per-second
(GOPS) $\Xi$ divided by the computational efficiency $\eta$ (in GOPS/W) of the
chosen FPGA/ASIC implementation: \todo{no channel estimation cost?}
\begin{align}
\label{eq:enc}
P_{\mathrm{enc}} &= P_{\mathrm{enc,s}}
+ \frac{f_{s,\mathrm{I}}}{\eta_{\mathrm{enc}}}\frac{14}{24}\frac{R_{\mathrm{DL}}}{\tilde B}, \\
P_{\mathrm{dec}} &= P_{\mathrm{dec,s}}
+ \frac{f_{s,\mathrm{I}}}{\eta_{\mathrm{dec}}}\frac{175}{6}\frac{R_{\mathrm{UL}}}{\tilde B},\\
\label{eq:pre}
P_{\mathrm{pre}} &= \bar M_{\mathrm{RF}}P_{s}
+ \frac{f_{s,\mathrm{I}}}{\eta_{\mathrm{pre}}}M_{\mathrm{RF}}\Xi_{\mathrm{pre}},\\
 P_{\mathrm{comb}} &= \bar M_{\mathrm{RF}}P_{s}
+ \frac{f_{s,\mathrm{I}}}{\eta_{\mathrm{comb}}}M_{\mathrm{RF}}\,2K,\\
\label{eq:ifft}
P_{\mathrm{IFFT}} &= P_{\mathrm{FFT}} = P_{\mathrm{IFFT,s}}
+ \frac{f_{s,\mathrm{II}}}{\eta_{\mathrm{IFFT}}}M_{\mathrm{RF}}\tfrac{3}{2}\log_2 Q_{\mathrm{IFFT}}, & & \\
\label{eq:dpd}
P_{\mathrm{DPD}} &= P_{\mathrm{DPD,s}}
+ \frac{f_{s,\mathrm{II}}}{\eta_{\mathrm{DPD}}}M_{\mathrm{RF}}\Xi_{\mathrm{DPD}}, & & \\
\label{eq:bb}
P_{\mathrm{filt,BB}} &= \bar M_{\mathrm{RF}}P_{s}
+ \frac{f_{s,\mathrm{II}}}{\eta_{\mathrm{BB}}}M_{\mathrm{RF}}\Xi_{\mathrm{BB}}, & &
\end{align}
where $R_{\mathrm{DL}},R_{\mathrm{UL}}$ are the delivered DL/UL sum rates,
$\tilde B$ the effective bandwidth,
$\bar M_{\mathrm{RF}}=\lceil M_{\mathrm{RF}}/32\rceil$ the number of FPGAs (one
per $32$ chains, sharing the precoder/combiner/BB-filter static power $P_{s}$),
and $Q_{\mathrm{IFFT}}$ the IFFT/FFT size. The two sampling rates are the
constellation rate $f_{s,\mathrm{I}}=\mu B$ (for encoding/decoding and
MIMO processing) and the oversampled rate
$f_{s,\mathrm{II}}=Q_{\mathrm{IFFT}}\Delta f$ (for the OFDM, DPD, and BB
operations), with $\mu\in[0,1]$ and subcarrier spacing $\Delta f$. For the
zero-forcing MIMO processor exploiting channel reciprocity, the precoder GOPS
is \todo{this changes for cell free depending on the precoder used}
\begin{equation}
\label{eq:gops_pre}
\Xi_{\mathrm{pre}} = 2K
+ \frac{1}{\upsilon_{\mathrm{coh}}}\!\left(\frac{K^3}{3M_{\mathrm{RF}}}
+ 3K^2 + K\right),
\end{equation}
with $\upsilon_{\mathrm{coh}}$ the number of samples per channel coherence
block over which the precoder is reused; the combiner only applies the
precoding matrix, giving $\Xi_{\mathrm{comb}}=2K$. The active digital
consumption per direction is the sum of its subcomponents,
\begin{align}
P_{\mathrm{dig,DL}} &= P_{\mathrm{enc}} + P_{\mathrm{pre}}
+ P_{\mathrm{IFFT}} + P_{\mathrm{DPD}} + P_{\mathrm{filt,BB}},\\
P_{\mathrm{dig,UL}} &= P_{\mathrm{dec}} + P_{\mathrm{comb}}
+ P_{\mathrm{IFFT}} + P_{\mathrm{filt,BB}},
\end{align}
and the frame-averaged digital consumption is
\begin{equation}
\label{eq:dig_avg}
\overline{P_{\mathrm{dig}}} =
\mathcal{A}_{\mathrm{DL}}^{\mathrm{dig}}\!\left(P_{\mathrm{dig,DL}}\right)
+ \mathcal{A}_{\mathrm{UL}}^{\mathrm{dig}}\!\left(P_{\mathrm{dig,UL}}\right).
\end{equation}

\subsection{Delivered rate and energy efficiency}
\label{sec:rate}
The achievable ergodic sum rate of direction $i$, which drives the
encoder/decoder consumption in \eqref{eq:enc}, follows from the per-user
SINRs under zero-forcing beamforming as
\begin{equation}
\label{eq:rate}
R_i = \bar x_i\,\tau_i\,(1-\tau_{i,\mathrm{sig}})\,\frac{\tilde B}{q_i}\,
\mathbb{E}\!\left\{\sum_{\nu=1}^{q_i}\sum_{k=1}^{K}
\log_2\!\left(1+\mathrm{SINR}_{i,k}[\nu]\right)\right\},
\end{equation}
where the pre-log collects the load, duration, signalling-overhead, and
effective-bandwidth-per-subcarrier factors. The base-station energy efficiency
then follows as the delivered rate per consumed watt,
\begin{equation}
\label{eq:ee}
\mathrm{EE} = \frac{R_{\mathrm{DL}} + R_{\mathrm{UL}}}{P_{\mathrm{cons}}}
\quad[\mathrm{bit/J}].
\end{equation}

\section{Distributed Massive MIMO System Model}

\subsection{Downlink System Model}

We consider the downlink of a distributed (cell-free) massive \gls{mimo}
network in which $L$ \glspl{ap}, each equipped with $M$ antennas, jointly serve
$K$ single-antenna users over the same time-frequency resources. The APs are
connected through fronthaul links to a central processing unit (CPU) that
encodes the data and computes the precoders, so that the $LM$ antennas act as a
single distributed array. Transmission uses OFDM with $Q$ data subcarriers,
indexed by $q=1,\dots,Q$; the model below is stated per subcarrier and the same
processing is applied on each of them.

Let $\vect{h}_{kl}[q]\inC{M}$ denote the channel from AP $l$ to user $k$ on
subcarrier $q$. Stacking the contributions of all APs gives the collective
channel
\begin{equation}
\label{eq:collective-channel}
\vect{h}_{k}[q] =
\begin{bmatrix} \vect{h}_{k1}[q] \\ \vdots \\ \vect{h}_{kL}[q] \end{bmatrix}
\inC{LM}.
\end{equation}
The channel realizations are drawn from the 3GPP TR 38.901 urban-microcell
(UMi) geometry-based stochastic channel model \cite{tr38901}, using the
implementation provided by the Sionna link-level simulator
\cite{hoydis2022sionna}. Each AP is modelled as a 38.901 base station carrying
an $M$-element half-wavelength \gls{ula} of omnidirectional elements and each
user as a single-antenna terminal. Every AP-user link is assigned a line-of-sight
or non-line-of-sight condition according to the distance-dependent LOS
probability of the UMi scenario, and its distance-dependent path loss, log-normal
shadow fading, and small-scale multipath then determine $\vect{h}_{kl}[q]$.

Because the model is specified directly in terms of channel realizations rather
than through a closed-form spatial correlation matrix, the large-scale fading
coefficient $\beta_{kl}$ between AP $l$ and user $k$, which drives the downlink
power control of \Cref{ssec:designs}, is obtained from the realizations as
the average per-antenna channel gain,
\begin{equation}
\label{eq:beta-estimate}
\beta_{kl} = \frac{1}{M}\,\expt{\norm{\vect{h}_{kl}[q]}^{2}},
\end{equation}
which is estimated by averaging $\abs{[\vect{h}_{kl}[q]]_{m}}^{2}$ over the $M$
antennas of AP $l$ and the $Q$ subcarriers. It captures the combined effect of
path loss and shadow fading (and, on a LOS link, of the specular component) and
so summarizes the slow, position-dependent part of the channel that the power
allocation acts on.

The CPU maps the data symbol $s_{k}[q]\in\mathbb{C}$ of user $k$, with unit
power $\expt{\abs{s_{k}[q]}^2}=1$ and mutually independent across users, onto the
$M$ antennas of each AP. AP $l$ transmits
\begin{equation}
\label{eq:ap-transmit}
\vect{x}_{l}[q] = \sum_{k=1}^{K} \vect{w}_{kl}[q]\, s_{k}[q]
= \mat{W}_{l}[q]\,\vect{s}[q],
\end{equation}
where $\vect{w}_{kl}[q]\inC{M}$ is the precoding vector that AP $l$ applies to
user $k$, $\vect{s}[q]=[\,s_{1}[q]\ \cdots\ s_{K}[q]\,]^{\intercal}\inC{K}$
collects the data symbols, and
$\mat{W}_{l}[q]=[\,\vect{w}_{1l}[q]\ \cdots\ \vect{w}_{Kl}[q]\,]\inC{M\times K}$
is the precoding matrix of AP $l$. Collecting the per-AP precoders of a given
user into
$\vect{w}_{k}[q]=[\,\vect{w}_{k1}[q]^{\intercal}\ \cdots\ \vect{w}_{kL}[q]^{\intercal}\,]^{\intercal}\inC{LM}$
yields the network-wide precoding matrix
\begin{equation}
\label{eq:network-precoder}
\mat{W}[q] = \big[\,\vect{w}_{1}[q]\ \cdots\ \vect{w}_{K}[q]\,\big]
\inC{LM\times K},
\end{equation}
whose $l$-th block of $M$ rows equals $\mat{W}_{l}[q]$, so that the stacked
transmit vector is
$\vect{x}[q]=[\,\vect{x}_{1}[q]^{\intercal}\ \cdots\ \vect{x}_{L}[q]^{\intercal}\,]^{\intercal}=\mat{W}[q]\,\vect{s}[q]\inC{LM}$.

The signal received by user $k$ on subcarrier $q$ is
$y_{k}[q]=\vect{h}_{k}[q]^{H}\vect{x}[q]+n_{k}[q]$, which, using
\eqref{eq:network-precoder}, separates into the desired term, the multi-user
interference, and the noise,
\begin{equation}
\label{eq:downlink-received}
y_{k}[q] =
\underbrace{\vect{h}_{k}[q]^{H}\vect{w}_{k}[q]\, s_{k}[q]}_{\text{desired signal}}
+ \underbrace{\sum_{\substack{i=1 \\ i\neq k}}^{K}
\vect{h}_{k}[q]^{H}\vect{w}_{i}[q]\, s_{i}[q]}_{\text{multi-user interference}}
+ \underbrace{\vphantom{\sum_{i=1}^{K}} n_{k}[q]}_{\text{noise}},
\end{equation}
where $n_{k}[q]\sim\cn{0}{\sigma^{2}}$ is the receiver noise with power
$\sigma^{2}$, independent across users and subcarriers and independent of the
channels and data. The scalar $\vect{h}_{k}[q]^{H}\vect{w}_{k}[q]$ is the
effective channel seen by user $k$: the distributed precoder turns the
$LM$-antenna channel into an equivalent single-antenna link.

Treating the multi-user interference in \eqref{eq:downlink-received} as noise,
the effective downlink \gls{sinr} of user $k$ on subcarrier $q$ is
\begin{equation}
\label{eq:downlink-sinr}
\mathrm{SINR}_{k}^{\mathrm{dl}}[q] =
\frac{\abs{\vect{h}_{k}[q]^{H}\vect{w}_{k}[q]}^{2}}
{\displaystyle\sum_{\substack{i=1 \\ i\neq k}}^{K}
\abs{\vect{h}_{k}[q]^{H}\vect{w}_{i}[q]}^{2} + \sigma^{2}} .
\end{equation}
With the effective channel known at the receiver, an achievable rate for user
$k$ on subcarrier $q$ is $\log_{2}\!\big(1+\mathrm{SINR}_{k}^{\mathrm{dl}}[q]\big)$,
and the delivered downlink sum rate $R_{\mathrm{DL}}$ follows by inserting these
per-user, per-subcarrier SINRs into the ergodic rate expression of
\Cref{sec:rate} (\eqref{eq:rate}), which in turn drives the encoder and MIMO
processing consumption of the power model.


\subsection{Precoder and Power Control Designs}
\label{ssec:designs}

Unlike co-located massive MIMO, where a single sum-power budget applies to the
whole array, each AP is a separate transmitter and is subject to its own power
constraint. Since the data symbols are independent and unit-power, the average
power radiated by AP $l$ across the OFDM block is
\begin{align}
\label{eq:power-constraint}
\sum_{q=1}^{Q}\expt{\norm{\vect{x}_{l}[q]}^{2}}
&= \sum_{q=1}^{Q}\sum_{k=1}^{K}\expt{\norm{\vect{w}_{kl}[q]}^{2}}
\leq \rho_{\max}, \\
  l&=1,\dots,L,
\end{align}
where $\rho_{\max}$ is the maximum downlink transmit power per AP, assumed equal
for all APs, and the expectation is taken over the data and channel
realizations.

It is convenient to split each precoder into a transmit direction and a power
control coefficient,
\begin{equation}
\label{eq:precoder-decomposition}
\vect{w}_{k}[q] = \sqrt{\rho_{k}}\,
\frac{\bar{\vect{w}}_{k}[q]}{\sqrt{\expt{\norm{\bar{\vect{w}}_{k}[q]}^{2}}}},
\end{equation}
where $\bar{\vect{w}}_{k}[q]\inC{LM}$ is an arbitrarily scaled vector that fixes
the spatial direction of the transmission to user $k$ and $\rho_{k}\geq 0$ is the
downlink power allocated to that user, so that
$\expt{\norm{\vect{w}_{k}[q]}^{2}}=\rho_{k}$. This separation isolates two design
tasks, the choice of the precoding directions $\{\bar{\vect{w}}_{k}[q]\}$ and the
power allocation that ties them to the per-AP budget in
\eqref{eq:power-constraint}. How both are computed depends on the level of
cooperation between the APs, so the two operation modes are treated in turn:
centralized operation (\Cref{sssec:centralized}), where the CPU designs all
precoders jointly and assigns one power coefficient per user, and local
operation (\Cref{sssec:local}), where each AP designs its own precoder and sets
its powers from local information alone. In both cases the allocation depends
mainly on the large-scale fading of \eqref{eq:beta-estimate} and can be held
fixed over many coherence blocks, whereas the directions are recomputed every
block.



\subsubsection{Centralized Operation}
\label{sssec:centralized}

In centralized operation the APs forward their received uplink pilots to the CPU,
which estimates the collective channels $\{\vect{h}_{k}[q]\}$ and designs all
precoders jointly \cite{demir2021foundations}. Exploiting TDD reciprocity, these
uplink estimates serve as the downlink \gls{csi}, so no downlink pilots are
needed. Because the CPU sees the whole $LM$-antenna array, it can make the
serving APs cancel one another's interference, using up to $LM$ spatial degrees
of freedom per user. Collecting the user channels into
$\mat{H}[q]=[\,\vect{h}_{1}[q]\ \cdots\ \vect{h}_{K}[q]\,]\inC{LM\times K}$, the
\gls{zf} precoder picks the directions that null all inter-user interference,
\begin{equation}
\label{eq:zf-precoder}
\bar{\mat{W}}[q] = \mat{H}[q]\big(\mat{H}[q]^{H}\mat{H}[q]\big)^{-1}
\inC{LM\times K},
\end{equation}
whose column $\bar{\vect{w}}_{k}[q]$ is orthogonal to every other user's channel
when the \gls{csi} is accurate; this requires $K\leq LM$ so that the Gram matrix
$\mat{H}[q]^{H}\mat{H}[q]$ is invertible. Zero-forcing disregards the noise, so
at low SNR it is improved by adding a diagonal loading $\lambda$ to the Gram
matrix, giving the regularized family
\begin{equation}
\label{eq:centralized-rzf}
\bar{\mat{W}}[q] = \mat{H}[q]\big(\mat{H}[q]^{H}\mat{H}[q]+\lambda\mat{I}_{K}\big)^{-1},
\end{equation}
which trades a little residual interference for robustness. With the loading set
to the noise power, $\lambda=\sigma^{2}$, this is the downlink \gls{mmse}
precoder, the dual of \gls{mmse} combining; under the perfect-\gls{csi}
assumption used here the \gls{mmse} precoder coincides with \gls{rzf} at that
loading, and the two differ only once a channel-estimation-error covariance is
modelled. Zero-forcing, the $\lambda=0$ limit, is the scheme adopted in the power
model of \Cref{sec:pwr_model}.

Because a centralized precoder is designed jointly across the APs, a single power
coefficient $\rho_{k}\geq 0$ is attached to each user and is common to all APs;
the direction $\bar{\vect{w}}_{k}[q]$ already fixes how that power is split over
the $LM$ antennas, and the direction is scaled to radiate $\rho_{k}$ as in
\eqref{eq:precoder-decomposition}. Two nominal allocations set $\{\rho_{k}\}$,
both spending the total network budget $P_{\mathrm{tot}}=L\rho_{\max}$. Equal
power gives every user the same share, $\rho_{k}=P_{\mathrm{tot}}/K$, whereas
fractional power control weights users by their aggregate large-scale gain
$\beta_{k}=\sum_{l=1}^{L}\beta_{kl}$,
\begin{equation}
\label{eq:centralized-fractional}
\rho_{k} = P_{\mathrm{tot}}\,
\frac{\beta_{k}^{\,v}}{\sum_{i=1}^{K}\beta_{i}^{\,v}},
\end{equation}
with the exponent $v$ trading throughput ($v>0$, favouring strong
users) against fairness ($v<0$, protecting weak users). These nominal powers
make the per-AP loads average to $\rho_{\max}$ but do not by themselves respect
each per-AP constraint in \eqref{eq:power-constraint}. Since the joint directions
cannot be rescaled per AP without breaking the interference nulling, a single
common factor
\begin{equation}
\label{eq:power-normalization}
\gamma = \sqrt{\frac{\rho_{\max}}{\max_{l}\tilde{p}_{l}}},
\qquad
\tilde{p}_{l} = \sum_{q=1}^{Q}\sum_{k=1}^{K}
\expt{\norm{\vect{w}_{kl}[q]}^{2}},
\end{equation}
rescales all precoders, where $\tilde{p}_{l}$ is the power AP $l$ would radiate
under the nominal allocation. The busiest AP then transmits at exactly
$\rho_{\max}$ and the others stay below their budget, while the common scalar
preserves the relative amplitudes and phases across APs and hence the
interference nulling of \eqref{eq:zf-precoder}. Because $\sum_{k}\rho_{k}=L\rho_{\max}$,
the per-AP powers average to $\rho_{\max}$, so $\gamma\leq 1$ and the constraint
is always feasible.

\subsubsection{Local Operation}
\label{sssec:local}

In local (distributed) operation the CPU only encodes the data, and each AP forms
its own precoding directions from its locally estimated channels
$\hat{\vect{h}}_{kl}[q]$, without sharing instantaneous \gls{csi}. An AP can then
suppress only the interference it itself generates, using the $M$ spatial degrees
of freedom of its own array, so the coherent cross-AP cancellation of the
centralized case is lost. The simplest choice is maximum-ratio transmission
(\gls{mr}), which points each beam along the local channel,
\begin{equation}
\label{eq:mr-precoder}
\bar{\vect{w}}_{kl}[q] = \vect{h}_{kl}[q],
\end{equation}
maximizing the signal power delivered to the desired user but ignoring the
interference it causes. Interference is instead suppressed locally by the
counterpart of zero-forcing applied to each AP's own users. Writing
$\mat{E}_{l}[q]=[\,\vect{h}_{1l}[q]\ \cdots\ \vect{h}_{Kl}[q]\,]\inC{M\times K}$
for AP $l$'s local channel matrix, the per-AP regularized precoder is
\begin{align}
\label{eq:local-rzf}
\mat{W}_{l}[q] &= \big(\mat{E}_{l}[q]\mat{E}_{l}[q]^{H}+\lambda\mat{I}_{M}\big)^{-1}
\mat{E}_{l}[q] \inC{M\times K},\\
 l&=1,\dots,L,
\end{align}
whose column $k$ is the local direction $\bar{\vect{w}}_{kl}[q]$, and the
collective direction $\bar{\vect{w}}_{k}[q]$ stacks the $L$ per-AP blocks. Local
zero-forcing is the $\lambda=0$ limit, which nulls the interference an AP causes
to its own users only when $M\geq K$; because each AP is small (typically
$M<K$), the loading is needed to keep the $M\times M$ inverse well conditioned,
giving local \gls{rzf} for a heuristic $\lambda$ and local \gls{mmse} for
$\lambda=\sigma^{2}$. As in the centralized case, under perfect \gls{csi} local
\gls{mmse} coincides with local \gls{rzf} at that loading. The $M\times M$
antenna-domain form is used because each AP's array is small.

Because each AP designs and normalizes its own block independently, the power is
parametrized per user per AP: AP $l$ normalizes its power to respect $\rho_{kl}\geq 0$ for user $k$ this is achieved as
\begin{align}
    \vect{w}_{kl}[q]=\sqrt{\rho_{kl}}\,\bar{\vect{w}}_{kl}[q]/\sqrt{\expt{\norm{\bar{\vect{w}}_{kl}[q]}^{2}}}.
\end{align}
Fractional power control is used, each AP splits its own budget
$\rho_{\max}$ among its users according to their local large-scale gains,
\begin{equation}
\label{eq:fractional-power}
\rho_{kl} = \rho_{\max}\,
\frac{\beta_{kl}^{\,v}}{\sum_{i=1}^{K}\beta_{il}^{\,v}},
\qquad l=1,\dots,L,
\end{equation}
so that $\sum_{k=1}^{K}\rho_{kl}=\rho_{\max}$ by construction and every AP meets
its per-array budget in \eqref{eq:power-constraint} with equality, using only
local quantities and without any CPU-side global rescaling. The exponent
$v$ plays the same throughput-versus-fairness role as above, with
$v=1/2$ recovering the heuristic of \cite{demir2021foundations}. Beyond
such statistics-based rules the powers could be optimized for a network utility
such as max-min fairness or sum rate, but those problems do not scale to very
large networks, so the heuristics above are the practical default.

\section{Distributed Massive MIMO Uplink System Model}
\label{sec:dmimo_ul_sysmodel}
We now consider the uplink of the same distributed (cell-free) massive \gls{mimo}
network as in \Cref{sec:dmimo_dl_sysmodel}. The $K$ users now transmit simultaneously on the same
time-frequency resources, each AP observes a superposition of all of them on its
$M$ antennas, and the per-AP observations are combined into one soft estimate per
user. Under \gls{tdd} operation the two directions occupy the same band within
one coherence block, so the propagation channels are the same $\vect{h}_{kl}[q]$
as in \eqref{eq:collective-channel} and the large-scale fading coefficients
$\beta_{kl}$ of \eqref{eq:beta-estimate} are unchanged. The downlink model
applies the array response as $\vect{h}_{k}[q]^{H}$, so reciprocity makes
$\vect{h}_{k}[q]$ itself the channel that carries the signal of user $k$ to the
$LM$ distributed antennas. 

User $k$ transmits $x_{k}[q]=\sqrt{p_{k}[q]}\,s_{k}[q]$ on subcarrier $q$, where
$s_{k}[q]\in\mathbb{C}$ is its unit-power data symbol,
$\expt{\abs{s_{k}[q]}^{2}}=1$, mutually independent across users, and
$p_{k}[q]\geq 0$ is the power it spends on that subcarrier. Each user is its own
transmitter and is subject to its own budget over the OFDM block,
\begin{equation}
\label{eq:ul-power-constraint}
\sum_{q=1}^{Q} p_{k}[q] = p_{k} \leq p_{\max},
\qquad k=1,\dots,K,
\end{equation}
with $p_{\max}$ the maximum transmit power of a user terminal. Since the power control below acts on the large-scale fading, which is flat over the band, the power is spread uniformly over the
subcarriers, $p_{k}[q]=p_{k}/Q$.

The signal received by AP $l$ on subcarrier $q$ is
\begin{equation}
\label{eq:ul-ap-received}
\vect{y}_{l}[q] = \sum_{k=1}^{K} \sqrt{p_{k}[q]}\,\vect{h}_{kl}[q]\, s_{k}[q]
+ \vect{n}_{l}[q] \inC{M},
\end{equation}
where $\vect{n}_{l}[q]\sim\cn{\vect{0}}{\sigma^{2}\mat{I}_{M}}$ is the receiver
noise, independent across APs and subcarriers and independent of the channels and
data. The noise power $\sigma^{2}$ is the same quantity as in
\eqref{eq:downlink-received}, but it is now referred to the AP receive chain, so
it is the one set by the noise figure of the LNAs and analog front end of
\Cref{sec:ana}. Stacking the $L$ observations gives the collective received vector
\begin{equation}
\label{eq:ul-collective-received}
\vect{y}[q] =
\begin{bmatrix} \vect{y}_{1}[q] \\ \vdots \\ \vect{y}_{L}[q] \end{bmatrix}
= \mat{H}[q]\,\mat{P}[q]^{1/2}\,\vect{s}[q] + \vect{n}[q] \inC{LM},
\end{equation}
with $\mat{H}[q]=[\,\vect{h}_{1}[q]\ \cdots\ \vect{h}_{K}[q]\,]\inC{LM\times K}$
the collective channel matrix already used in \Cref{sssec:centralized},
$\vect{s}[q]=[\,s_{1}[q]\ \cdots\ s_{K}[q]\,]^{\intercal}\inC{K}$,
$\vect{n}[q]\sim\cn{\vect{0}}{\sigma^{2}\mat{I}_{LM}}$, and
$\mat{P}[q]^{1/2}\inC{K\times K}$ the diagonal matrix with
$[\mat{P}[q]^{1/2}]_{k,k}=\sqrt{p_{k}[q]}$ and zero off-diagonal entries.

The data symbol of user $k$ is estimated by applying a receive combining vector
$\vect{v}_{k}[q]\inC{LM}$ to \eqref{eq:ul-collective-received}, which separates
into the desired term, the multi-user interference, and the noise,
\begin{align}
\label{eq:ul-estimate}
\hat{s}_{k}[q] ={}& \vect{v}_{k}[q]^{H}\vect{y}[q] \nonumber\\
={}& \underbrace{\sqrt{p_{k}[q]}\,\vect{v}_{k}[q]^{H}\vect{h}_{k}[q]\, s_{k}[q]}
_{\text{desired signal}} \nonumber\\
&+ \underbrace{\sum_{\substack{i=1 \\ i\neq k}}^{K}
\sqrt{p_{i}[q]}\,\vect{v}_{k}[q]^{H}\vect{h}_{i}[q]\, s_{i}[q]}
_{\text{multi-user interference}}
+ \underbrace{\vphantom{\sum_{i=1}^{K}}\vect{v}_{k}[q]^{H}\vect{n}[q]}
_{\text{noise}} .
\end{align}
Writing $\vect{v}_{k}[q]=[\,\vect{v}_{k1}[q]^{\intercal}\ \cdots\
\vect{v}_{kL}[q]^{\intercal}\,]^{\intercal}$ in blocks of $M$ entries, with
$\vect{v}_{kl}[q]\inC{M}$ the part applied to AP $l$, the estimate is the sum of
per-AP contributions,
\begin{equation}
\label{eq:ul-distributive}
\hat{s}_{k}[q] = \sum_{l=1}^{L} \vect{v}_{kl}[q]^{H}\vect{y}_{l}[q].
\end{equation}
Linear combining is therefore distributive over the APs, and this is the property
the functional splits of \Cref{ssec:split} exploit: an AP that holds its own
combining coefficients can form its partial sum locally and forward one scalar
stream per user instead of $M$ sample streams.

Treating the multi-user interference in \eqref{eq:ul-estimate} as noise, and
using that the noise term has variance
$\sigma^{2}\norm{\vect{v}_{k}[q]}^{2}$, the effective uplink \gls{sinr} of user
$k$ on subcarrier $q$ is
\begin{equation}
\label{eq:ul-sinr}
\mathrm{SINR}_{k}^{\mathrm{ul}}[q] =
\frac{p_{k}[q]\,\abs{\vect{v}_{k}[q]^{H}\vect{h}_{k}[q]}^{2}}
{\displaystyle\sum_{\substack{i=1 \\ i\neq k}}^{K}
p_{i}[q]\,\abs{\vect{v}_{k}[q]^{H}\vect{h}_{i}[q]}^{2}
+ \sigma^{2}\norm{\vect{v}_{k}[q]}^{2}} .
\end{equation}
With the effective channel known at the decoder, an achievable rate for user $k$ on subcarrier $q$ is
$\log_{2}\!\big(1+\mathrm{SINR}_{k}^{\mathrm{ul}}[q]\big)$, and inserting these
per-user, per-subcarrier SINRs into the ergodic rate expression of
\Cref{sec:rate} \eqref{eq:rate} gives the delivered per-user uplink rate
\begin{align}
\label{eq:ul-rate-user}
R_{\mathrm{UL},k} ={}& \bar x_{\mathrm{UL}}\,\tau_{\mathrm{UL}}\,
(1-\tau_{\mathrm{UL,sig}})\,\frac{\tilde B}{Q} \nonumber\\
&\times \expt{\sum_{q=1}^{Q}\log_{2}\!\left(1
+\mathrm{SINR}_{k}^{\mathrm{ul}}[q]\right)},
\end{align}
with $R_{\mathrm{UL}}=\sum_{k=1}^{K}R_{\mathrm{UL},k}$ the delivered uplink sum
rate. It drives the channel-decoder consumption of \Cref{sec:dig} and, through
\eqref{eq:fh_ul}, the traffic-dependent part of the fronthaul.


\subsection{Combiner and Power Control Designs}
\label{ssec:ul_designs}

As in the downlink, the design splits into the choice of the receive directions
$\{\vect{v}_{k}[q]\}$, which depends on how much the APs cooperate, and the
choice of the transmit powers $\{p_{k}\}$, which is a per-user allocation driven
by the large-scale fading. The two are treated in turn, and the directions are
recomputed every coherence block whereas the powers can be held fixed over many
blocks.

Throughout, and consistently with \Cref{ssec:designs}, we assume \emph{perfect}
\gls{csi}: every combiner is built from the true channels,
$\hat{\vect{h}}_{k}[q]=\vect{h}_{k}[q]$, so there is no estimation error and no
pilot contamination. The pilot overhead is nevertheless charged to the rate
through the frame parameters of \Cref{sec:frame}. In the downlink-only
evaluation of \Cref{ssec:designs} the uplink phase carried nothing but pilots,
$\tau_{\mathrm{UL,sig}}=1$, which leaves \eqref{eq:ul-rate-user} at zero. An
uplink that carries data requires $\tau_{\mathrm{UL,sig}}<1$, and the coherence
block then has to be split explicitly: with $\tau_{c}$ samples per block of which
$\tau_{p}$ are pilots, the signalling fraction of the frame is
$\tau_{\mathrm{UL}}\,\tau_{\mathrm{UL,sig}}=\tau_{p}/\tau_{c}$ and the uplink
data prelog of \eqref{eq:ul-rate-user} is
$\tau_{\mathrm{UL}}(1-\tau_{\mathrm{UL,sig}})=\tau_{\mathrm{UL}}-\tau_{p}/\tau_{c}$,
which requires $\tau_{p}/\tau_{c}\leq\tau_{\mathrm{UL}}$. Mutually orthogonal
pilots additionally require $\tau_{p}\geq K$; this is what makes the
no-contamination assumption above self-consistent, and it is the assumption under
which the perfect-\gls{csi} model is the idealisation of a real pilot phase rather
than a different system.


\subsubsection{Centralized Operation}
\label{sssec:ul_centralized}

In centralized operation the APs forward their received samples to the \gls{cpu},
which holds $\vect{y}[q]$ in full and combines over all $LM$ antennas jointly
\cite{demir2021foundations}. This is the highest level of cooperation, the
receive counterpart of the joint precoding of \Cref{sssec:centralized}, and it
corresponds to splits S1 and S2 of \Cref{ssec:split}. Minimizing the mean squared
error $\expt{\abs{s_{k}[q]-\hat{s}_{k}[q]}^{2}}$ over $\vect{v}_{k}[q]$ for the
model \eqref{eq:ul-collective-received} gives the \gls{mmse} combiner
\begin{equation}
\label{eq:ul-mmse-combiner}
\vect{v}_{k}[q] = \left(\mat{H}[q]\,\mat{P}[q]\,\mat{H}[q]^{H}
+ \sigma^{2}\mat{I}_{LM}\right)^{-1}\vect{h}_{k}[q],
\end{equation}
which requires inverting an $LM\times LM$ matrix. Applying the push-through
identity $(\mat{A}\mat{B}+\sigma^{2}\mat{I})^{-1}\mat{A}
=\mat{A}(\mat{B}\mat{A}+\sigma^{2}\mat{I})^{-1}$ with $\mat{A}=\mat{H}[q]\mat{P}[q]$
and $\mat{B}=\mat{H}[q]^{H}$, and discarding the diagonal factor that the scale
invariance of \eqref{eq:ul-sinr} makes irrelevant, the same directions are
obtained from a $K\times K$ inverse,
\begin{equation}
\label{eq:ul-centralized-rzf}
\mat{V}[q] = \mat{H}[q]\left(\mat{H}[q]^{H}\mat{H}[q]
+ \sigma^{2}\mat{P}[q]^{-1}\right)^{-1} \inC{LM\times K},
\end{equation}
whose column $k$ is $\vect{v}_{k}[q]$ up to a scalar. This is the uplink form of
the regularized family \eqref{eq:centralized-rzf}, with the heuristic loading
$\lambda\mat{I}_{K}$ replaced by the physically determined
$\sigma^{2}\mat{P}[q]^{-1}$; for equal powers $p_{k}[q]=p$ the two coincide at
$\lambda=\sigma^{2}/p$, which is the sense in which the downlink \gls{rzf}
precoder is the dual of \gls{mmse} combining. Letting the loading vanish, which
is the high-SNR limit $\sigma^{2}/p\to 0$, gives the \gls{zf} combiner
\begin{equation}
\label{eq:ul-zf-combiner}
\mat{V}[q] = \mat{H}[q]\big(\mat{H}[q]^{H}\mat{H}[q]\big)^{-1},
\end{equation}
for which $\vect{v}_{k}[q]^{H}\vect{h}_{i}[q]=0$ for all $i\neq k$, so all
inter-user interference is removed and \eqref{eq:ul-sinr} reduces to
$p_{k}[q]/(\sigma^{2}\norm{\vect{v}_{k}[q]}^{2})$. As in the downlink this needs
$K\leq LM$ for the Gram matrix to be invertible, and it disregards the noise, so
it is the scheme of choice only when the loading is negligible. Zero-forcing is
the scheme adopted in the power model of \Cref{sec:pwr_model} and its distributed
extension.

Because the combiner is scale-invariant, the centralized design carries no power
allocation of its own; the transmit powers are set at the users by
\eqref{eq:ul-fractional-power} below. The \gls{cpu} does need the powers
$\{p_{k}\}$ to build \eqref{eq:ul-centralized-rzf}, but these are large-scale
quantities refreshed once per large-scale fading realisation, not once per
coherence block.


\subsubsection{Local Operation and Large-Scale Fading Decoding}
\label{sssec:ul_local}

In local operation each AP combines its own observation with its own \gls{csi}
and forwards one scalar per user, so no instantaneous \gls{csi} crosses the
fronthaul. This is split S3 of \Cref{ssec:split}. AP $l$ forms
$\hat{s}_{kl}[q]=\vect{v}_{kl}[q]^{H}\vect{y}_{l}[q]$ using either \gls{mr}
combining, which points along the local channel,
\begin{equation}
\label{eq:ul-mr-combiner}
\vect{v}_{kl}[q] = \vect{h}_{kl}[q],
\end{equation}
or the local regularized combiner built from AP $l$'s own channel matrix
$\mat{E}_{l}[q]=[\,\vect{h}_{1l}[q]\ \cdots\ \vect{h}_{Kl}[q]\,]\inC{M\times K}$,
\begin{align}
\label{eq:ul-local-rzf}
\mat{V}_{l}[q] &= \big(\mat{E}_{l}[q]\mat{P}[q]\mat{E}_{l}[q]^{H}
+\sigma^{2}\mat{I}_{M}\big)^{-1}\mat{E}_{l}[q] \inC{M\times K},\\
 l&=1,\dots,L,
\end{align}
whose column $k$ is $\vect{v}_{kl}[q]$. This is \eqref{eq:local-rzf} with the
heuristic loading again replaced by its physical value, and it inverts an
$M\times M$ matrix because the per-AP array is small. An AP can suppress only the
interference it observes itself, using the $M$ degrees of freedom of its own
array, so the coherent cross-AP interference rejection of
\eqref{eq:ul-centralized-rzf} is lost; the loading is what keeps the inverse well
conditioned in the usual regime $M<K$.

What remains is how the \gls{cpu} fuses the $L$ local estimates. It applies one
weight per AP and user,
\begin{equation}
\label{eq:ul-lsfd}
\hat{s}_{k}[q] = \sum_{l=1}^{L} a_{kl}[q]^{*}\,\hat{s}_{kl}[q],
\end{equation}
collected into $\vect{a}_{k}[q]=[\,a_{k1}[q]\ \cdots\ a_{kL}[q]\,]^{\intercal}
\inC{L}$. By \eqref{eq:ul-distributive} this is exactly the collective combiner
$\vect{v}_{k}[q]$ whose $l$-th block is $a_{kl}[q]\vect{v}_{kl}[q]$, so
\eqref{eq:ul-sinr} continues to apply unchanged and the fusion rule is part of the
combiner rather than a separate stage. Two choices are used. Equal weighting,
$a_{kl}[q]=1$, simply adds the local estimates and needs nothing at the \gls{cpu}
beyond the streams themselves. \Gls{lsfd} instead optimizes the weights from the
channel statistics: collecting the per-AP effective gains
$g_{kil}[q]\triangleq\vect{v}_{kl}[q]^{H}\vect{h}_{il}[q]$ into
$\vect{g}_{ki}[q]=[\,g_{ki1}[q]\ \cdots\ g_{kiL}[q]\,]^{\intercal}\inC{L}$, and
writing $\mat{D}_{k}[q]$ for the diagonal matrix with entries
$\norm{\vect{v}_{kl}[q]}^{2}$ (diagonal because the noise is independent across
APs), the weights that maximize the resulting use-and-then-forget \gls{sinr} are
\begin{align}
\label{eq:ul-lsfd-weights}
\vect{a}_{k}[q] ={}& \Big(\sum_{i=1}^{K} p_{i}[q]\,
\expt{\vect{g}_{ki}[q]\vect{g}_{ki}[q]^{H}}
+ \sigma^{2}\expt{\mat{D}_{k}[q]}\Big)^{-1} \nonumber\\
&\times \sqrt{p_{k}[q]}\ \expt{\vect{g}_{kk}[q]},
\end{align}
where the expectations are over the small-scale fading for fixed user positions
\cite{demir2021foundations}.
\todo{verify \eqref{eq:ul-lsfd-weights} against the LSFD expression in
\cite{demir2021foundations} before it is used in the evaluation.}
These weights are statistical, not instantaneous, so they are computed once per
large-scale fading realisation and cost the fronthaul nothing per coherence
block. The three fusion rules therefore span three cooperation levels at the same
per-AP combiner: local combining at a single AP, equal-weight summation at the
\gls{cpu}, and statistically weighted summation. Only the centralized combiner of
\eqref{eq:ul-centralized-rzf} uses instantaneous \gls{csi} at the \gls{cpu}, and
the SE expressions of the different levels must not be mixed.

Two things should be flagged about how \eqref{eq:ul-sinr} is used at these lower
levels. When the fusion weights are statistical, the \gls{cpu} does not know the
instantaneous effective channel $\vect{v}_{k}[q]^{H}\vect{h}_{k}[q]$, and the
rigorous achievable rate is then the use-and-then-forget bound, in which the mean
effective channel is treated as the useful gain and its fluctuation as extra
interference. Evaluating \eqref{eq:ul-rate-user} with the instantaneous
\eqref{eq:ul-sinr}, as is done in the downlink of \Cref{sec:dmimo_dl_sysmodel},
assumes instead that the effective channel is known wherever the decoding
happens, and is therefore optimistic for local operation by the amount of
residual channel hardening.
\todo{decide whether the local-operation rate is reported with the instantaneous
SINR, for consistency with the downlink, or with the use-and-then-forget bound,
which is the achievable one; the gap grows as $M$ shrinks.}


\subsubsection{Uplink Power Control}
\label{sssec:ul_power}

The powers $\{p_{k}\}$ are set at the users from the large-scale fading and, by
\eqref{eq:ul-power-constraint}, are constrained one user at a time. Full power,
$p_{k}=p_{\max}$, is the natural reference, but it is the worst case for
fairness: a user close to an AP then swamps a distant one sharing the same
resources. Fractional power control weights the users by their aggregate
large-scale gain $\beta_{k}=\sum_{l=1}^{L}\beta_{kl}$,
\begin{equation}
\label{eq:ul-fractional-power}
p_{k} = p_{\max}\,
\frac{\beta_{k}^{\,v}}{\max_{i}\beta_{i}^{\,v}},
\qquad k=1,\dots,K,
\end{equation}
which satisfies $p_{k}\leq p_{\max}$ for either sign of the exponent $v$, since
the normalizing maximum is attained by the strongest user when $v>0$ and by the
weakest when $v<0$. The exponent plays the same throughput-versus-fairness role
as in \eqref{eq:centralized-fractional}, with $v=0$ recovering full power and
$v=-1$ inverting the channel statistically, so that every user arrives at the
network with the same average total gain $p_{k}\beta_{k}$. The normalization by a
maximum rather than by a sum is what differs from the downlink rule
\eqref{eq:centralized-fractional}: there a shared budget had to be divided among
the users, whereas here each user has its own budget and the only role of the
denominator is to keep the strongest user within it. As a result the uplink has
no counterpart of the common rescaling \eqref{eq:power-normalization}, and no
feasibility question arises.

The uplink transmit power is spent at the users and therefore does not appear in
the network consumption $P_{\mathrm{net}}$ of \eqref{eq:pnet}, which counts the
APs, the fronthaul, and the \gls{cpu} only. The asymmetry this creates in the
power model is worth stating plainly: in the downlink phase the APs pay for the
PAs of \Cref{ssec:dmimo_pa}, whose consumption scales with the radiated power,
whereas in the uplink phase they pay only for the receive front end
\eqref{eq:ana_ul_ap} and the digital processing \eqref{eq:dig_ap}, both of which
are essentially independent of $p_{\max}$. Uplink power control therefore acts on
the delivered rate \eqref{eq:ul-rate-user} and on nothing else in
\eqref{eq:ee_net}, which makes it a pure spectral-efficiency lever rather than an
energy one, and it is the reason the uplink and downlink phases of the same frame
have very different consumption profiles.


\subsection{Complexity of combiners}
\todo{compute complexity of the combiners so this can be used in the power model;
under TDD reciprocity the centralized combiner reuses the matrix already
inverted for \eqref{eq:gops_cen}, so only the application cost $2K'N$ should be
charged in the uplink, whereas \gls{lsfd} adds the per-user $L\times L$ solve of
\eqref{eq:ul-lsfd-weights} once per large-scale fading realisation.}

\section{D-MIMO Power Model}
\label{sec:dmimo_pwr}

The model of \Cref{sec:pwr_model} describes one central site: a single co-located array
whose antennas, converters, and processors share a physical location, a clock, and a power
supply. The distributed system of \Cref{sec:dmimo_sysmodel} replaces that site by
$L$ \glspl{ap} of $M$ antennas each, spread over the coverage area and tied to a
\gls{cpu} by fronthaul links. Three things change as a result. First, the hardware
blocks become per-AP quantities that must be summed over the network, each at its
own operating point. Second, the transport between the \glspl{ap} and the \gls{cpu} has its own power consumption, wich is not modeled in the single-site model. Third, the digital
processing chain is split across two locations (\gls{ap} and \gls{cpu}). This requires design choices about where a given operation runs. This section extends the parametric model along those three axes while keeping the frame-averaging machinery of \Cref{sec:frame} intact, so that the single-site model is recovered as a special case by setting $L=1$ with no fronthaul and no \gls{cpu}.

Throughout, $l=1,\dots,L$ indexes APs, $k=1,\dots,K$ users, $q=1,\dots,Q$ subcarriers, and $i\in\{\mathrm{DL},\mathrm{UL}\}$ the link direction, matching \Cref{sec:dmimo_sysmodel}. Hardware parameters carried over from \Cref{sec:pwr_model} keep their symbols and now refer to a single AP.

\subsection{Network power decomposition}
\label{ssec:dmimo_structure}

Let $\mathcal{S}\subseteq\{1,\dots,L\}$ be the set of APs that are switched on;
an AP outside $\mathcal{S}$ is in deep sleep and draws a fraction
$\delta^{\mathrm{ds}}\in[0,1]$ of its base consumption
$P^{0}_{\mathrm{AP},l}$. The total distributed network consumption is
\begin{align}
\label{eq:pnet}
P_{\mathrm{net}} ={}& \sum_{l\in\mathcal{S}}
\left(P_{\mathrm{AP},l} + \frac{\overline{P_{\mathrm{FH},l}}}{\eta_{\mathrm{FH,sc}}}\right)
\nonumber\\
&+ \sum_{l\notin\mathcal{S}} \delta^{\mathrm{ds}}
\left(P^{0}_{\mathrm{AP},l} + \frac{P^{0}_{\mathrm{FH},l}}{\eta_{\mathrm{FH,sc}}}\right)
+ \frac{\overline{P_{\mathrm{CPU}}}}{\eta_{\mathrm{CPU,sc}}},
\end{align}
where the consumption of an active AP retains the three-block structure of
\eqref{eq:pcons},
\begin{equation}
\label{eq:pap}
P_{\mathrm{AP},l} =
\frac{\overline{P_{\mathrm{dig},l}}}{\eta_{\mathrm{dig,sc}}}
+ \frac{\overline{P_{\mathrm{ana},l}}}{\eta_{\mathrm{ana,sc}}}
+ \frac{\overline{P_{\mathrm{PA},l}}}{\eta_{\mathrm{PA,sc}}},
\end{equation}
and $\eta_{\mathrm{FH,sc}},\eta_{\mathrm{CPU,sc}}\in(0,1]$ are the
supply-and-cooling efficiencies of the fronthaul transceivers and of the central
unit. The CPU efficiency plays the role of an inverse power-usage effectiveness:
the central unit sits in a room with its own cooling, which is not the same
environment as an outdoor AP enclosure, so $\eta_{\mathrm{CPU,sc}}$ should not be
set equal to the $\eta_{c,\mathrm{sc}}$ values fitted for a macro base station. \todo{find proper values, look at power usage effectiveness (PUE) of datacenters}
Setting $L=1$, $\mathcal{S}=\{1\}$, $M=M_{\mathrm{ant}}$, and dropping the
fronthaul and CPU terms returns \eqref{eq:pcons} unchanged.

The set $\mathcal{S}$ determines how many \glspl{ap} serve the network. As a first step we assume that all $L$ \glspl{ap} are active. 

\todo{see if we want to keep the next paragraph:}
The set $\mathcal{S}$ is the first genuinely new degree of freedom. A co-located
array can shut down individual chains but not the site; a distributed deployment
can shut down whole APs while the network keeps serving, and the coverage
redundancy of cell-free operation is what makes that possible. AP selection
therefore acts directly on \eqref{eq:pnet}, and \Cref{rem:fh_ceiling} shows that
it also relieves the fronthaul.

\subsection{Serving clusters, per-AP load, and the shared frame}
\label{ssec:dmimo_frame}

In the fully distributed operation of \Cref{sec:dmimo_sysmodel} every AP serves
every user. The power model should not assume this, because restricting each AP
to a subset of users is the standard way of keeping cell-free operation scalable
and is also the main lever on fronthaul traffic. Let
$\mathcal{D}_{l}\subseteq\{1,\dots,K\}$ be the set of users served by AP $l$ and
$\mathcal{L}_{k}=\{l:k\in\mathcal{D}_{l}\}$ the serving cluster of user $k$, with
$K_{l}\triangleq\abs{\mathcal{D}_{l}}$ and $\abs{\mathcal{L}_{k}}$ its
cardinality. Full cell-free operation is the special case
$\mathcal{D}_{l}=\{1,\dots,K\}$ for all $l$, so $K_{l}=K$ and
$\abs{\mathcal{L}_{k}}=L$. A user outside $\mathcal{D}_{l}$ has
$\vect{w}_{kl}[q]=\vect{0}$, which leaves the signal model of
\eqref{eq:ap-transmit} untouched.

The frame timing, in contrast, is shared. Under \gls{tdd} the APs must switch
between directions together, since an AP still transmitting while a neighbour
receives creates cross-link interference between nodes that are far closer to
each other than to any user. The frame parameters
$\tau_{\mathrm{DL}},\tau_{\mathrm{UL}},\tau_{i,\mathrm{sig}}$ of
\Cref{sec:frame} are therefore network-wide constants, and only the load is
allowed to differ between APs. Repeating the construction of
\eqref{eq:load_avg} over the users AP $l$ actually serves gives the per-AP
average physical resource load
\begin{equation}
\label{eq:load_ap}
\bar x_{i,l} = \frac{N_{a,i}}{N_{i}}\,\frac{K_{i,l}}{K_{\max}},
\end{equation}
with $K_{i,l}=\abs{\mathcal{D}_{l}}$ counted over the users active in direction
$i$ and $K_{\max}$ the network-wide number of spatial streams. The
frame-averaging function \eqref{eq:frameavg} is reused with this per \gls{ap} load,
\begin{align}
\label{eq:frameavg_ap}
\mathcal{A}_{i,l}^{c}\!\left(P^{\mathrm a},P^{\mathrm s},P^{0}\right)
\triangleq{}&
\bar x_{i,l}\,\tau_i\,(1-\tau_{i,\mathrm{sig}})\,P^{\mathrm a}
+ \tau_i\,\tau_{i,\mathrm{sig}}\,P^{\mathrm s} \nonumber\\
&+ (1-\bar x_{i,l})\,\tau_i\,(1-\tau_{i,\mathrm{sig}})\,\delta^{\mu}_c\,P^{0}
\nonumber\\
&+ (1-\tau_i)\,\delta^{0}_c\,P^{0},
\end{align}
and the shorthand $\mathcal{A}_{i,l}^{c}(P)\triangleq
\mathcal{A}_{i,l}^{c}(P,P,P)$ is kept for the blocks whose consumption does not
depend on the operating mode. Under full cell-free operation all APs share the
same load, $\bar x_{i,l}=\bar x_{i}$, and \eqref{eq:frameavg_ap} collapses to
\eqref{eq:frameavg}.

\paragraph{Reconciling the two overhead accountings}
\todo{double check this part and work out some inconcistencies in the notation used above wrt prolog factors}
\Cref{sec:dmimo_sysmodel} charges the pilot overhead through the prelog
$(1-\tau_{p}/\tau_{c})$ of the coherence-block bookkeeping, while
\Cref{sec:rate} charges it through $\tau_{i}(1-\tau_{i,\mathrm{sig}})$. These are
two descriptions of the same time budget and must be identified rather than
multiplied. Under TDD reciprocity the pilots are sent in the uplink, so the
uplink signalling phase of \Cref{sec:frame} \emph{is} the pilot phase of the
coherence block, and writing $\tau_{d}^{\mathrm{DL}}$ for the number of samples
per block that carry downlink data, the identification is
\begin{equation}
\label{eq:overhead_map}
\tau_{\mathrm{UL}}\,\tau_{\mathrm{UL,sig}} = \frac{\tau_{p}}{\tau_{c}},
\qquad
\tau_{\mathrm{DL}}\,(1-\tau_{\mathrm{DL,sig}}) = \frac{\tau_{d}^{\mathrm{DL}}}{\tau_{c}},
\end{equation}
with $\tau_{d}^{\mathrm{DL}}\leq\tau_{c}-\tau_{p}$, the slack being the uplink
data and the downlink demodulation reference signals. The downlink-only prelog
$(1-\tau_{p}/\tau_{c})$ used in \Cref{sec:dmimo_sysmodel} is the special case
$\tau_{d}^{\mathrm{DL}}=\tau_{c}-\tau_{p}$, in which the whole block after the
pilots carries downlink data and neither uplink data nor downlink reference
signals are scheduled. The rate expression \eqref{eq:rate} then applies the
prelog once and must not apply it again. A second bookkeeping detail: the
subcarrier count is written $q_{i}$ in \eqref{eq:rate}, $Q_{\max}$ in
\eqref{eq:load_inst}, and $Q$ in \Cref{sec:dmimo_sysmodel}, while $q$ is also the
subcarrier index. We use $Q$ for the number of evaluated data subcarriers
throughout this section. \todo{align the symbols for the subcarrier count across
\Cref{sec:pwr_model} and \Cref{sec:dmimo_sysmodel}.}

\subsection{Per-AP power amplifiers}
\label{ssec:dmimo_pa}

Each AP is an independent transmitter with its own budget
\eqref{eq:power-constraint}, so the single quantity $P_{T,\mathrm{DL}}$ of
\Cref{sec:pa} becomes a vector of per-AP transmit powers. From the precoders of
\Cref{ssec:designs}, the power actually radiated by AP $l$ over the OFDM block is
\begin{equation}
\label{eq:ptl}
P_{T,l} \triangleq \sum_{q=1}^{Q}\sum_{k=1}^{K}
\expt{\norm{\vect{w}_{kl}[q]}^{2}} \leq \rho_{\max},
\end{equation}
which the two operation modes fill in differently. Local operation meets the
constraint with equality by construction, $P_{T,l}=\rho_{\max}$ for every AP,
because \eqref{eq:fractional-power} splits each AP's own budget among its users.
Centralized operation cannot rescale APs independently without breaking the
interference nulling, so the common factor of \eqref{eq:power-normalization}
gives $P_{T,l}=\gamma^{2}\tilde p_{l}$ with $\gamma^{2}=\rho_{\max}/\max_{l'}\tilde
p_{l'}$, and only the busiest AP reaches its budget. It is convenient to work
with the per-AP power utilisation
\begin{equation}
\label{eq:util}
u_{l} \triangleq \frac{P_{T,l}}{\rho_{\max}} \in [0,1].
\end{equation}

As was done in the centralized power model we assume that the per-AP power is shared equally across its antennas, the per-PA
output power is $P_{a,l}=P_{T,l}/M$ and the frame-averaged PA consumption of AP
$l$ follows \eqref{eq:pa_avg} with the per-AP load,
\begin{equation}
\label{eq:pa_avg_ap}
\overline{P_{\mathrm{PA},l}} = M\,
\mathcal{A}_{\mathrm{DL},l}^{\mathrm{PA}}\!\left(
P_{\mathrm{PA}}(P_{a,l}),\,
P_{\mathrm{PA}}(\zeta_{\mathrm{sig}}P_{\max}),\,
P_{\mathrm{PA}}(0)\right),
\end{equation}
with $P_{\max}$ the maximum PA output power and $\zeta_{\mathrm{sig}}$ the
fraction of it used during signalling, both as in \eqref{eq:pa_avg}.

Note that the equal-split assumption is not exact: the per-antenna powers
$\sum_{q}\sum_{k}\expt{\abs{[\vect{w}_{kl}[q]]_{m}}^{2}}$ differ across
$m=1,\dots,M$ for any precoder that is not isotropic. Because the efficiency
exponent $\alpha\in[0,1]$ of \eqref{eq:pa} makes $p\mapsto p^{\alpha}$ concave,
and hence $P_{\mathrm{PA}}(\cdot)$ concave, Jensen's inequality gives
$\frac{1}{M}\sum_{m}P_{\mathrm{PA}}(P_{a,lm}) \leq P_{\mathrm{PA}}(P_{a,l})$: the
equal-split assumption is a conservative upper bound on the true PA consumption,
and it is tight when the precoder loads the antennas of an AP uniformly.

\begin{remark}[PA sizing and the static floor]
\label{rem:pa_sizing}
The relation between $P_{\max}$ and the AP budget must be stated, because it
decides how the PA static term scales with the deployment. If each PA is
dimensioned for its share of the AP budget, $P_{\max}=\rho_{\max}/M$, then
substituting \eqref{eq:pa} into \eqref{eq:pa_avg_ap} makes the data-mode PA
consumption of AP $l$ independent of $M$,
\begin{equation}
\label{eq:pa_compact}
M\,P_{\mathrm{PA}}(P_{a,l}) =
\frac{\rho_{\max}}{\eta_{\mathrm{PA,max}}}
\Big[\xi + (1-\xi)\,u_{l}^{\alpha}\Big],
\end{equation}
where $\xi\in[0,1]$ is the static-to-dynamic consumption ratio and
$\eta_{\mathrm{PA,max}}$ the maximum PA efficiency of \eqref{eq:pa}. A fully
loaded AP therefore consumes exactly $\rho_{\max}/\eta_{\mathrm{PA,max}}$, and
the network static floor is $L\,\xi\rho_{\max}/\eta_{\mathrm{PA,max}}$, growing
with the number of APs rather than with the number of antennas. If instead
$P_{\max}$ is a fixed component rating independent of $M$, the floor is
$LM\,\xi P_{\max}/\eta_{\mathrm{PA,max}}$ and grows with the total antenna count.
The two conventions differ by a factor $M P_{\max}/\rho_{\max}$ and must not be
mixed. The three PA parameters $\xi$, $\alpha$, and $\eta_{\mathrm{PA,max}}$
were fitted for macro-cell amplifiers and should be refitted for the small
amplifiers of a distributed AP rather than reused.
\todo{source $\xi,\alpha,\eta_{\mathrm{PA,max}}$ for AP-class PAs.}
\end{remark}

\begin{remark}[Cost of the per-AP constraint]
\label{rem:gamma}
\todo{check this}
Under centralized operation with either nominal allocation of
\Cref{sssec:centralized}, the nominal powers satisfy $\sum_{l}\tilde
p_{l}=\sum_{k}\rho_{k}=L\rho_{\max}$, so the average utilisation is
\begin{equation}
\label{eq:ubar}
\bar u \triangleq \frac{1}{L}\sum_{l=1}^{L} u_{l}
= \frac{\gamma^{2}}{L\rho_{\max}}\sum_{l=1}^{L}\tilde p_{l}
= \gamma^{2}.
\end{equation}
The common rescaling factor of \eqref{eq:power-normalization} is therefore
exactly the fraction of the network power budget that is actually radiated.
Combining \eqref{eq:pa_compact} with \eqref{eq:ubar} and the concavity of
$u\mapsto u^{\alpha}$ bounds the network PA consumption in the data mode from
below by
\begin{equation}
\label{eq:pa_network_bound}
\sum_{l=1}^{L} M\,P_{\mathrm{PA}}(P_{a,l}) \geq
\frac{L\rho_{\max}}{\eta_{\mathrm{PA,max}}}
\Big[\xi + (1-\xi)\,\gamma^{2\alpha}\Big],
\end{equation}
with equality when all APs are equally loaded. The per-AP constraint thus costs
twice: the network radiates only a fraction $\gamma^{2}$ of $L\rho_{\max}$, yet
the static floor $L\xi\rho_{\max}/\eta_{\mathrm{PA,max}}$ is paid in full
whenever the PAs are not asleep. An imbalanced deployment, in which one AP is
much busier than the rest and drives $\gamma$ down, is penalised on both counts.
\end{remark}

\subsection{Per-AP analog front end}
\label{ssec:dmimo_ana}

A distributed AP carries few antennas, so hybrid beamforming has nothing to
amortise and the arrays are fully digital: $M_{\mathrm{RF}}=M$ and
$M_{\mathrm{PS}}=0$. Specialising \eqref{eq:ana_dl} and \eqref{eq:ana_ul}
accordingly, the working-mode analog consumption of one AP is
\begin{align}
\label{eq:ana_dl_ap}
P_{\mathrm{ana,DL},l} &= P_{\mathrm{LO}} + M\big(
2P_{\mathrm{DAC}} + 2P_{\mathrm{filt,RF}} + 2P_{\mathrm{mix}}\big),\\
\label{eq:ana_ul_ap}
P_{\mathrm{ana,UL},l} &= P_{\mathrm{LO}} + M\big(
2P_{\mathrm{ADC}} + 2P_{\mathrm{filt,RF}} + 2P_{\mathrm{mix}}
+ P_{\mathrm{LNA}}\big),
\end{align}
and the frame-averaged analog block of AP $l$ is
\begin{align}
\label{eq:ana_avg_ap}
\overline{P_{\mathrm{ana},l}} ={}&
\mathcal{A}_{\mathrm{DL},l}^{\mathrm{ana}}\!\left(P_{\mathrm{ana,DL},l}\right)
+ \mathcal{A}_{\mathrm{UL},l}^{\mathrm{ana}}\!\left(P_{\mathrm{ana,UL},l}\right)
\nonumber\\
&+ P_{\mathrm{ana,misc}} + P_{\mathrm{sync}} .
\end{align}
The local oscillator is the first term to change character. In the co-located
model a single LO is shared by $M_{\mathrm{ant}}$ chains and contributes a
constant; here every AP needs its own, so the network pays $L\,P_{\mathrm{LO}}$
and the term grows with the deployment even though it is small per node.

The term $P_{\mathrm{sync}}$ has no counterpart in \Cref{sec:ana} and is added
here. Coherent joint transmission across physically separate nodes requires the
APs to share a frequency and phase reference and requires the TDD chains to be
reciprocity-calibrated; the precoder of \eqref{eq:zf-precoder} nulls
interference only if the relative phases across APs are the ones the CPU assumed.
Whether the reference is distributed over the fronthaul, over a dedicated timing
network, or over the air, it costs power at every AP continuously, which is why
$P_{\mathrm{sync}}$ sits outside the direction-dependent averaging. It is a real
cost of distributing the array and one that a co-located design does not pay.
\todo{source a value for $P_{\mathrm{sync}}$; over-the-air calibration also
consumes frame time, which is not yet charged to $\tau_{i,\mathrm{sig}}$.}

\subsection{Functional split of Digital Signal Processing}
\label{ssec:split}

In the single-site model every digital operation runs at the base station. In a
distributed deployment the encoder, the precoder computation, and the precoder
application can each sit at the CPU or at the AP, and the choice sets both the
digital consumption and the fronthaul rate. We distinguish three splits, which
span the practical range and map onto the operation modes of
\Cref{ssec:designs}; \Cref{tab:split} summarises where each operation ends up.

\begin{table}[t]
\centering
\caption{Node at which each digital operation runs, per functional split. The
uplink combiner follows the same allocation as the downlink precoder. Splits S1
and S2 both realise the centralized precoding of \Cref{sssec:centralized} and
deliver the same rate; S3 is local operation (\Cref{sssec:local}) and delivers
the lower rate of local precoding. $^{\dagger}$Listed for completeness but not
priced, since \Cref{sec:dmimo_sysmodel} assumes perfect CSI \todo{add this}.}
\label{tab:split}
\ra{1.15}
\begin{tabular}{@{}llccc@{}}
\toprule
Operation & Count & S1 & S2 & S3 \\
\midrule
Encoding / decoding    & \eqref{eq:enc}  & CPU & CPU & CPU \\
Channel estimation$^{\dagger}$ & ---     & CPU & CPU & AP  \\
Precoder computation   & \eqref{eq:gops_cen}, \eqref{eq:gops_loc}
                                         & CPU & CPU & AP  \\
Precoder application   & $2K'N$          & CPU & AP  & AP  \\
OFDM (I)FFT            & \eqref{eq:ifft} & AP  & AP  & AP  \\
Predistortion          & \eqref{eq:dpd}  & AP  & AP  & AP  \\
Baseband filter        & \eqref{eq:bb}   & AP  & AP  & AP  \\
\bottomrule
\end{tabular}
\end{table}

Split~S1 forwards precoded samples. The CPU encodes, computes the precoder of
\eqref{eq:zf-precoder}, and applies it, sending AP $l$ the $M$ streams of
frequency-domain samples it must radiate; the AP performs only the OFDM
modulation, predistortion, baseband filtering, and conversion. 

Split~S2 forwards coefficients and data. The CPU encodes and computes $\mat{W}_{l}[q]$ but AP $l$
applies it locally, so the fronthaul carries the $K_{l}$ data streams plus the
precoding coefficients, refreshed once per coherence block. 


Split~S3 is local operation as defined in \Cref{sssec:local}: the CPU only encodes, and AP $l$
estimates its own channels and computes and applies the precoder of
\eqref{eq:local-rzf} from local information, so the fronthaul carries data alone.
Splits~S1 and~S2 both realise the centralized precoding of
\Cref{sssec:centralized} and therefore deliver the same rate; they differ only in
where the work is done and what crosses the fronthaul. Split~S3 delivers the
lower rate of local precoding.

\subsection{Digital signal processing at the APs and the CPU}
\label{ssec:dmimo_dig}

The per-chain digital costs of \Cref{sec:dig} carry over unchanged, evaluated
with $M_{\mathrm{RF}}=M$ at each AP. Every AP performs the OFDM modulation and
demodulation, the predistortion of its own PAs, and the baseband filtering
whatever the split, so
\begin{align}
\label{eq:dig_ap}
P_{\mathrm{dig,DL},l} &= P_{\mathrm{IFFT}} + P_{\mathrm{DPD}}
+ P_{\mathrm{filt,BB}} + P_{\mathrm{MIMO,DL},l},\\
P_{\mathrm{dig,UL},l} &= P_{\mathrm{IFFT}} + P_{\mathrm{filt,BB}}
+ P_{\mathrm{MIMO,UL},l},
\end{align}
with $\overline{P_{\mathrm{dig},l}}=
\mathcal{A}_{\mathrm{DL},l}^{\mathrm{dig}}(P_{\mathrm{dig,DL},l})+
\mathcal{A}_{\mathrm{UL},l}^{\mathrm{dig}}(P_{\mathrm{dig,UL},l})$. Only the
MIMO-processing terms $P_{\mathrm{MIMO},i,l}$ depend on the split, and they are
where the architecture actually shows.

As in \Cref{sec:dig}, a MIMO-processing term is a static part plus
$f_{s,\mathrm{I}}\Xi/\eta_{\mathrm{pre}}$, with $\Xi$ its complex operations per
sample. Two operations must be kept apart, because the split moves them
independently. Applying an already-computed precoder to $K'$ users over $N$
chains costs $2K'N$ operations per sample, the count already used for the
combiner in \Cref{sec:dig}. Computing it costs the Cholesky terms of
\eqref{eq:gops_pre}, amortised over the $\upsilon_{\mathrm{coh}}$ samples of a
coherence block. Under centralized operation the CPU inverts the $K\times K$ Gram
matrix of the collective channel $\mat{H}[q]\inC{LM\times K}$, so
\eqref{eq:gops_pre} applies with $M_{\mathrm{RF}}$ replaced by the total number
of distributed antennas $LM$,
\begin{equation}
\label{eq:gops_cen}
\Xi_{\mathrm{pre}}^{\mathrm{cen}} =
\frac{1}{\upsilon_{\mathrm{coh}}}\left(\frac{K^{3}}{3}
+ 3K^{2}LM + K\,LM\right),
\end{equation}
written here as a network total rather than per chain. Under local operation AP
$l$ instead inverts the $M\times M$ matrix of \eqref{eq:local-rzf}, which is the
same Cholesky-based accounting with the roles of the antenna and user dimensions
exchanged,
\begin{equation}
\label{eq:gops_loc}
\Xi_{\mathrm{pre},l}^{\mathrm{loc}} =
\frac{1}{\upsilon_{\mathrm{coh}}}\left(\frac{M^{3}}{3}
+ 3M^{2}K_{l} + M\,K_{l}\right).
\end{equation}
Setting $L=1$ and $M=M_{\mathrm{RF}}$ in \eqref{eq:gops_cen} recovers the
computation part of $M_{\mathrm{RF}}\Xi_{\mathrm{pre}}$ in \eqref{eq:gops_pre}
exactly.

The three splits then distribute these two counts over the two nodes as
\begin{align}
\label{eq:gops_split}
\Xi^{\mathrm{CPU}}_{i} &= \iota_{\mathrm{S1}}\,2K\,LM
+ \iota_{\mathrm{S1,S2}}\,\iota_{\mathrm{DL}}\,\Xi_{\mathrm{pre}}^{\mathrm{cen}},\\
\Xi_{i,l}^{\mathrm{AP}} &= \iota_{\mathrm{S2,S3}}\,2K_{l}M
+ \iota_{\mathrm{S3}}\,\iota_{\mathrm{DL}}\,\Xi_{\mathrm{pre},l}^{\mathrm{loc}},
\end{align}
where each $\iota_{\bullet}\in\{0,1\}$ equals one for the split or direction
named in its subscript and zero otherwise, and the corresponding powers follow
\eqref{eq:pre},
\begin{equation}
\label{eq:pmimo}
P_{\mathrm{MIMO},i,l} = \bar M_{\mathrm{RF}}P_{s}
+ \frac{f_{s,\mathrm{I}}}{\eta_{\mathrm{pre}}}\,\Xi_{i,l}^{\mathrm{AP}},
\qquad
\bar M_{\mathrm{RF}} = \Big\lceil \tfrac{M}{32} \Big\rceil ,
\end{equation}
and likewise for $P_{\mathrm{MIMO},i}^{\mathrm{CPU}}$ with
$\Xi^{\mathrm{CPU}}_{i}$. As in \Cref{sec:dig}, the computation is charged to the
downlink only, since TDD reciprocity lets the same matrix serve the uplink
combiner, while the application is charged in both directions. Every operation
then appears exactly once across \eqref{eq:gops_split}, so no work is counted
twice: under S2, in particular, the AP is billed only for applying the precoder
and the CPU only for computing it. Because linear combining is distributive over
the APs, $\hat s_{k}=\sum_{l}\vect{v}_{kl}^{H}\vect{y}_{l}$, an AP holding its
combining coefficients can form its own partial sum locally, so S2 attains the
centralized uplink rate while forwarding $K_{l}$ scalar streams instead of $M$
sample streams.

Summed over the network, \eqref{eq:gops_loc} grows linearly in $L$ at fixed $M$
while \eqref{eq:gops_cen} grows linearly in $LM$ and additionally carries the
$K^{3}$ term. For the dense, small-$M$ deployments of interest the centralized
cost is dominated by $3K^{2}LM$ and the local cost by $3M^{2}K_{l}$, so
centralized precoding is more expensive by roughly a factor $K/M$ per antenna.

Channel encoding and decoding stay at the CPU in all three splits, since that is
where the payload enters and leaves the network. The CPU consumption is
\begin{align}
\label{eq:cpu}
\overline{P_{\mathrm{CPU}}} ={}& P_{\mathrm{CPU},0}
+ \mathcal{A}_{\mathrm{DL}}^{\mathrm{dig}}\!\left(P_{\mathrm{enc}}
+ P_{\mathrm{MIMO,DL}}^{\mathrm{CPU}}\right) \nonumber\\
&+ \mathcal{A}_{\mathrm{UL}}^{\mathrm{dig}}\!\left(P_{\mathrm{dec}}
+ P_{\mathrm{MIMO,UL}}^{\mathrm{CPU}}\right),
\end{align}
with $P_{\mathrm{enc}}$ and $P_{\mathrm{dec}}$ from \eqref{eq:enc} driven by the
network sum rates $R_{\mathrm{DL}}$ and $R_{\mathrm{UL}}$, and
$P_{\mathrm{CPU},0}$ the always-on consumption of the central unit. The averaging
in \eqref{eq:cpu} uses the network-wide load $\bar x_{i}$, not a per-AP one.

\begin{remark}[Loss of static-power amortisation]
\label{rem:fpga}
The precoder, combiner, and baseband filter of \Cref{sec:dig} share a static
power $P_{s}$ per FPGA, one per $\bar M_{\mathrm{RF}}=\lceil
M_{\mathrm{RF}}/32\rceil$ group of chains. A co-located array of
$M_{\mathrm{ant}}$ antennas amortises that static power over 32 chains at a time;
a distributed network of $L$ APs with $M<32$ antennas each pays it $L$ times,
once per AP, for a total of $LM$ chains. At $L=40$ and $M=4$ the network runs
$40$ FPGAs where a co-located array of the same $160$ antennas would run $5$, an
eightfold increase in this static term. Chopping a large array into many small
ones leaves the per-chain digital costs untouched but destroys the sharing of
the fixed costs, and this is a general feature rather than an artefact of the
value $32$.
\end{remark}

\subsection{Fronthaul consumption}
\label{ssec:dmimo_fh}

Following the cell-free power model of \cite{ngo2018total}, the fronthaul link
serving AP $l$ consumes a traffic-independent part and a part proportional to the
rate it carries,
\begin{equation}
\label{eq:fh_basic}
P_{\mathrm{FH},l} = P_{\mathrm{FH},0} + \Pi_{\mathrm{FH}}\,R_{\mathrm{FH},l},
\end{equation}
with $P_{\mathrm{FH},0}$ in watts, the traffic-dependent coefficient
$\Pi_{\mathrm{FH}}$ in W per bit/s, and $R_{\mathrm{FH},l}$ the fronthaul rate in
bit/s. The static part covers the optical transceivers at both ends and depends
on the link length and the topology; the traffic part covers the
throughput-dependent consumption of the transport chain. In the parametrisation
of \cite{ngo2018total} these are $P_{\mathrm{FH},0}=\SI{0.825}{\watt}$ per link
and $\Pi_{\mathrm{FH}}=\SI{0.25}{\watt}$ per Gbit/s.
\todo{these two values are taken from a sub-6\,GHz cell-free study; check that
they are representative of the fronthaul technology assumed here before using
them for FR3 numbers.}

\paragraph{Fronthaul rate per split}
What $R_{\mathrm{FH},l}$ contains is set by the functional split of
\Cref{ssec:split}. Let $b_{\mathrm{FH}}$ be the number of bits per real sample
on the fronthaul, so a complex sample costs $2b_{\mathrm{FH}}$ bits. Under S1 the
link carries $M$ streams of frequency-domain samples at the constellation rate
$f_{s,\mathrm{I}}$, in both directions,
\begin{equation}
\label{eq:fh_s1}
R_{\mathrm{FH},l}^{\mathrm{S1}} = 2\,b_{\mathrm{FH}}\,M\,f_{s,\mathrm{I}},
\end{equation}
which becomes $2b_{\mathrm{FH}}Mf_{s,\mathrm{II}}$ if the CPU also performs the
OFDM modulation and time-domain samples are forwarded. Under S2 the link carries
the payload of the served users plus the precoding coefficients, an $M\times
K_{l}$ complex matrix per subcarrier refreshed once per coherence block of
$\upsilon_{\mathrm{coh}}$ samples,
\begin{equation}
\label{eq:fh_s2}
R_{\mathrm{FH},l}^{\mathrm{S2,DL}} = \sum_{k\in\mathcal{D}_{l}} R_{\mathrm{DL},k}
+ \frac{f_{s,\mathrm{I}}}{\upsilon_{\mathrm{coh}}}\,2\,b_{\mathrm{FH}}\,M\,K_{l}.
\end{equation}
Under S3 only the payload crosses in the downlink,
$R_{\mathrm{FH},l}^{\mathrm{S3,DL}}=\sum_{k\in\mathcal{D}_{l}}R_{\mathrm{DL},k}$.
In the uplink both S2 and S3 forward $K_{l}$ streams of per-AP partial sums
rather than $M$ streams of raw samples,
\begin{equation}
\label{eq:fh_ul}
R_{\mathrm{FH},l}^{\mathrm{S2,UL}} = R_{\mathrm{FH},l}^{\mathrm{S3,UL}}
= 2\,b_{\mathrm{FH}}\,K_{l}\,f_{s,\mathrm{I}},
\end{equation}
which for S2 costs nothing in rate, by the distributivity argument of
\Cref{ssec:dmimo_dig}. The power-control coefficients are refreshed only once per
large-scale fading realisation and are neglected, as in \cite{ngo2018total}; the
precoding coefficients of \eqref{eq:fh_s2} are not, because they change every
coherence block.

The signalling mode deserves separate treatment. Centralized operation requires
the CPU to see the raw received pilots in order to estimate the collective
channel, so during the uplink signalling phase the link must carry samples even
under S2, giving $R^{\mathrm{sig}}_{\mathrm{FH},l}=2b_{\mathrm{FH}}M
f_{s,\mathrm{I}}$ for S1 and S2 alike. Under S3 the pilots are consumed locally
and $R^{\mathrm{sig}}_{\mathrm{FH},l}\approx 0$. Centralization therefore imposes
a sample-level fronthaul during pilot transmission whatever the data-phase split.

\paragraph{Frame averaging}
The link is bidirectional and stays up across the whole frame, so its static term
is charged once per link with a single duty cycle rather than once per direction.
Let
\begin{align}
\label{eq:fh_duty}
\theta_{l} \triangleq \sum_{i\in\{\mathrm{DL},\mathrm{UL}\}}
\tau_{i}\Big[\tau_{i,\mathrm{sig}}
+ (1-\tau_{i,\mathrm{sig}})\,\bar x_{i,l}\Big]
\end{align}
be the fraction of the frame in which the link is in a working mode. With
$\delta^{\mu}_{\mathrm{FH}}$ the reduction factor of an idle transceiver, the
frame-averaged fronthaul consumption is
\begin{equation}
\label{eq:fh_avg}
\overline{P_{\mathrm{FH},l}} =
\Big[\theta_{l} + (1-\theta_{l})\,\delta^{\mu}_{\mathrm{FH}}\Big]
P_{\mathrm{FH},0}
+ \Pi_{\mathrm{FH}}\,\overline{R}_{\mathrm{FH},l},
\end{equation}
with the frame-averaged fronthaul rate obtained by reusing
\eqref{eq:frameavg_ap} with a zero base level, since an idle link carries no
traffic,
\begin{equation}
\label{eq:fh_rate_avg}
\overline{R}_{\mathrm{FH},l} = \sum_{i\in\{\mathrm{DL},\mathrm{UL}\}}
\mathcal{A}_{i,l}^{\mathrm{FH}}\!\left(
R^{i}_{\mathrm{FH},l},\,R^{i,\mathrm{sig}}_{\mathrm{FH},l},\,0\right).
\end{equation}
Optical transceivers wake slowly and hold most of their consumption while idle,
so $\delta^{\mu}_{\mathrm{FH}}$ is close to one in practice and the static
fronthaul term is nearly a constant per active link. This is what makes
$\mathcal{S}$, rather than the duty cycle, the effective lever on the static part.

\begin{remark}[Where the frame structure actually enters]
\label{rem:no_double}
The map from rate to power in \eqref{eq:fh_basic} is linear, and
$\mathcal{A}_{i,l}^{\mathrm{FH}}$ with a zero base is itself linear in its
arguments. For the data-sharing splits S2 and S3 the quantity to insert in
\eqref{eq:fh_rate_avg} is therefore the \emph{peak} rate during the data mode,
and since \eqref{eq:rate} already delivers $R_{\mathrm{DL},k}$ with the prelog
$\bar x_{i}\tau_{i}(1-\tau_{i,\mathrm{sig}})$ applied, the downlink data
contribution collapses back to $\sum_{k\in\mathcal{D}_{l}}R_{\mathrm{DL},k}$.
Charging the delivered rate directly is thus equivalent to frame-averaging the
peak rate, and doing both would count the frame structure twice. For the
sample-forwarding split S1 the situation is the opposite: $R^{\mathrm{S1}}_{\mathrm{FH},l}$
in \eqref{eq:fh_s1} is a hardware constant with no dependence on the delivered
rate at all, and the frame structure is then the only mechanism that reduces the
fronthaul load. The frame structure and the traffic dependence are therefore not
two independent knobs; each split makes one of them the active one.
\end{remark}

\begin{remark}[Fronthaul ceiling on energy efficiency]
\label{rem:fh_ceiling}
Under a data-sharing split the same payload is duplicated to every AP of a user's
serving cluster, so the network fronthaul load is
$\sum_{l}\sum_{k\in\mathcal{D}_{l}}R_{\mathrm{DL},k}=
\sum_{k}\abs{\mathcal{L}_{k}}R_{\mathrm{DL},k}$. Keeping only this term in
\eqref{eq:pnet} and using $\sum_{k}\abs{\mathcal{L}_{k}}R_{\mathrm{DL},k}\geq
(\min_{k}\abs{\mathcal{L}_{k}})\sum_{k}R_{\mathrm{DL},k}$ bounds the downlink
energy efficiency by
\begin{equation}
\label{eq:ee_ceiling}
\mathrm{EE}_{\mathrm{DL}} \leq
\frac{1}{\Pi_{\mathrm{FH}}\,\min_{k}\abs{\mathcal{L}_{k}}},
\end{equation}
independently of the bandwidth, the transmit power, and the precoder. Full
cell-free operation has $\abs{\mathcal{L}_{k}}=L$, so with $L=40$ and
$\Pi_{\mathrm{FH}}=\SI{0.25}{\watt}$ per Gbit/s the network cannot exceed
$\SI{100}{\mega\bit\per\joule}$ however well everything else is designed.
Shrinking the serving clusters raises the ceiling in direct proportion, which is
why AP selection is the dominant lever on cell-free energy efficiency and why
this term, absent from the co-located model, can decide the comparison on its own.
\end{remark}

\subsection{Delivered rate and network energy efficiency}
\label{ssec:dmimo_ee}

The per-user downlink rate follows from the SINR of \eqref{eq:downlink-sinr}
inserted into the ergodic rate expression \eqref{eq:rate},
\begin{align}
\label{eq:rate_user}
R_{\mathrm{DL},k} ={}& \bar x_{\mathrm{DL}}\,\tau_{\mathrm{DL}}\,
(1-\tau_{\mathrm{DL,sig}})\,\frac{\tilde B}{Q} \nonumber\\
&\times \expt{\sum_{q=1}^{Q}\log_{2}\!\left(1
+\mathrm{SINR}_{k}^{\mathrm{dl}}[q]\right)},
\end{align}
with $R_{\mathrm{DL}}=\sum_{k=1}^{K}R_{\mathrm{DL},k}$ the delivered sum rate.
The per-user decomposition, which the co-located model did not need, is required
here because the fronthaul term \eqref{eq:fh_s2} charges each AP only for the
users in its cluster. The network energy efficiency is then
\begin{equation}
\label{eq:ee_net}
\mathrm{EE}_{\mathrm{net}} =
\frac{R_{\mathrm{DL}} + R_{\mathrm{UL}}}{P_{\mathrm{net}}}
\quad[\mathrm{bit/J}],
\end{equation}
which reduces to \eqref{eq:ee} for a single AP without fronthaul. Because a
distributed and a co-located deployment cover the same area with different
hardware counts, a comparison at equal $L M$ and equal total radiated power is
the meaningful one, and the area energy efficiency
$\mathrm{EE}_{\mathrm{net}}/A$ over the served area $A$ is the fairer figure when
the deployments differ in extent.

A capacity constraint should accompany \eqref{eq:ee_net}. If the fronthaul links
have a finite capacity $C_{\mathrm{FH}}$, then $R_{\mathrm{FH},l}\leq
C_{\mathrm{FH}}$ must hold for every $l$, and under S1 this is a hard constraint
on $M$ and the sampling rate that is independent of the traffic. When it binds,
the rate of \eqref{eq:rate_user} is not achievable as written and the model
returns an optimistic energy efficiency.

\subsection{What the extension does not yet capture}
\label{ssec:dmimo_gaps}

Several effects are left out and are flagged here rather than silently absorbed
into the parameters. Fronthaul quantization is modelled only through its rate:
the $b_{\mathrm{FH}}$ bits per sample of \eqref{eq:fh_s1} add a distortion term
to the effective SINR of \eqref{eq:downlink-sinr} that is not present in
\Cref{sec:dmimo_sysmodel}, so S1 is charged for its fronthaul load but not
debited for the accuracy it loses. Channel estimation processing is likewise
unpriced: \Cref{sec:dmimo_sysmodel} assumes perfect \gls{csi}, so neither the
$LM\times K$ estimation at the CPU under centralized operation nor the $M\times
K$ local estimation under S3 appears in \eqref{eq:cpu} or \eqref{eq:dig_ap},
although both scale with the deployment. The synchronization and
reciprocity-calibration overhead is charged as a continuous power
$P_{\mathrm{sync}}$ but not as frame time, even though over-the-air calibration
consumes symbols that would otherwise carry data. Finally, the reduction factors
$\delta^{\mu}_{c}$ and $\delta^{0}_{c}$, the supply-and-cooling efficiencies, and
the PA parameters are all inherited from a macro base station and are the most
likely source of systematic error when they are applied unchanged to a
small outdoor AP. \todo{refit the AP-class hardware parameters, or report results
as a sensitivity range over them.}



%\linespread{0.98}
\bibliographystyle{IEEEtran}
\bibliography{IEEEabrv,mybib}

\end{document}
